% concatenated from WO_and_Pi_1_3_Uniformization.tex
\documentclass[a4paper,11pt]{article}
\usepackage[utf8]{inputenc}

\usepackage{boxedminipage}
\usepackage{amsfonts}
\usepackage{amsmath} 
\usepackage{amssymb}
\usepackage{graphicx}
\usepackage{amsthm}
\usepackage{t1enc}
\usepackage{subfig}
\usepackage{cancel}
 \usepackage{textcmds}
 \usepackage{mathabx}

\newtheorem{theorem}{Theorem}[section]
\newtheorem{lemma}[theorem]{Lemma}
\newtheorem{proposition}[theorem]{Proposition}
\newtheorem{corollary}[theorem]{Corollary}
\newtheorem{definition}[theorem]{Definition}
\newtheorem{example}[theorem]{Example}
\newtheorem{fact}[theorem]{Fact}
\newtheorem{question}{Question}
\newtheorem*{theorem*}{Theorem}
\newtheorem{claim}{Claim}

\newcommand{\forceP}{\mathbb{P}}
\newcommand{\forceQ}{\mathbb{Q}}
\newcommand{\forceR}{\mathbb{R}}
\newcommand{\forceL}{\mathbb{L}}


\newcommand{\ZFC}{\mathsf{ZFC}}
\newcommand{\ZF}{\mathsf{ZF}}
\newcommand{\ZFP}{\mathsf{ZF}^-}
\newcommand{\AC}{\mathsf{AC}}
\newcommand{\PFA}{\mathsf{PFA}}
\newcommand{\BPFA}{\mathsf{BPFA}}
\newcommand{\MRP}{\mathsf{MRP}}
\newcommand{\MM}{\mathsf{MM}}
\newcommand{\CH}{\mathsf{CH}}
\newcommand{\GCH}{\mathsf{GCH}}
\newcommand{\MA}{\mathsf{MA}_{\omega_1}}
\newcommand{\PD}{\mathsf{PD}}


\def\undertilde#1{\mathord{\vtop{\ialign{##\crcr
$\hfil\displaystyle{#1}\hfil$\crcr\noalign{\kern1.5pt\nointerlineskip}
$\hfil\tilde{}\hfil$\crcr\noalign{\kern1.5pt}}}}}



%opening
\title{A Universe with a $\Delta^1_n$-definable well-order of the reals, $\CH$ and $\Pi^1_n$-Uniformization}
\author{ Stefan Hoffelner$^{1}$\footnote{This research was funded in whole by the Austrian Science Fund (FWF) Grant-DOI 10.55776/P37228. }  }
\date{
    $^1$TU Wien \\
    \today }
\begin{document}

\maketitle

\begin{abstract}
This paper constructs a universe where $\Pi^1_3$-uniformization is true, the Continuum Hypothesis holds yet it possesses a $\Delta^1_3$-definable well-order of its reals. The method can be lifted to canonical inner models with finitely many Woodin cardinals to produce universes of $\mathsf{CH}$, $\Pi^1_n$-uniformization and where additionally a $\Delta^1_n$-definable well-order of the reals exist.
\end{abstract}

\section{Introduction}

This paper explores the interaction between two important properties of the real numbers: the $\Pi^1_n$-uniformization property and the $\Delta^1_n$-definable well-orderings of the reals.

To recall, for a set \( A \subset \omega^\omega \times \omega^\omega \), we say that a function \( f \) is a uniformization (or uniformizing function) of \( A \) if \( f \) is a partial function \( f: \omega^\omega \to \omega^\omega \), with the domain of \( f \) being \( \text{pr}_1(A) \), and the graph of \( f \) being a subset of \( A \).

\begin{definition}
We say that a projective pointclass \( \Gamma \in \{ \Sigma^1_n, \Pi^1_n \mid n \in \omega \} \) has the uniformization property if every element of \( \Gamma \) has a uniformization that also belongs to \( \Gamma \).
\end{definition}

J. Addison observed in \cite{Addison} that a good  well-order of the reals, definable by a \( \Delta^1_n \)-formula, implies the \( \Sigma^1_m \)-uniformization property for every \( m \ge n \). A \( \Delta^1_n \)-definable well-order \( < \) of the reals is considered a \emph{good} \( \Delta^1_n \)-well-order if its order type is \( \omega_1 \), and the relation \( <_I \subset (\omega^\omega)^2 \) is defined as follows:
\[
x <_I y \Leftrightarrow \{(x)_n \mid n \in \omega\} = \{z \mid z < y\},
\]
where \( (x)_n \) represents a recursive partition of \( x \) into \( \omega \)-many reals, and this relation is \( \Delta^1_n \)-definable.

Novikov's result shows that uniformization cannot hold for both \( \Sigma^1_n \) and \( \Pi^1_n \) simultaneously for any \( n \in \omega \). Specifically, \( \Pi^1_n \)-uniformization and a good \( \Delta^1_n \)-definable well-order cannot coexist. It is natural to ask whether it is consistent with \( \mathsf{CH} \) that a \( \Delta^1_n \)-definable well-order of the reals can be combined with the \( \Pi^1_n \)-uniformization property. This paper demonstrates that it is indeed possible.

\begin{theorem*}
There exists a generic extension of \( L \) in which there is a \( \Delta^1_3 \)-definable well-order of the reals, \( \mathsf{CH} \) holds, and the \( \Pi^1_3 \)-uniformization property holds.
\end{theorem*}
The proof can be adapted to apply to canonical inner models with Woodin cardinals, 

\begin{theorem*}
Let \( M_n \) be the canonical inner model with \( n \) Woodin cardinals. Then there exists a generic extension of \( M_n \) in which there is a \( \Delta^1_{n+3} \)-definable well-order of the reals, \( \mathsf{CH} \) holds, and the \( \Pi^1_{n+3} \)-uniformization property holds.
\end{theorem*}
The techniques used to obtain such a universe are based on the work of \cite{Ho1}, \cite{Ho2}, and \cite{Ho3}, to which this paper owes much. However, we employ a different coding machine, resulting in a countable support iteration rather than a finite (or mixed support iteration, depending on the ground model). We believe that this new machine is of independent interest.

Put in a bigger context, this article contributes to the study of the behaviour of separation, reduction and uniformization on the projective hierarchy in the absence of projective determinacy (see \cite{Ho5} for a recent and very different result).



\section{Preliminaries}


The forcings used in this construction are well-known, but we will briefly introduce them and highlight their key properties.

\begin{definition}(see \cite{BHK})
For a stationary set \( R \subset \omega_1 \), the club-shooting forcing for \( R \), denoted \( \forceP_R \), consists of conditions \( p \) which are countable functions from \( \alpha+1 < \omega_1 \) to \( R \) whose image is a closed set. The partial order \( \forceP_R \) is ordered by end-extension.
\end{definition}

The club-shooting forcing \( \forceP_R \) serves as a classic example of an \( R \)-\emph{proper forcing}. A forcing \( \forceP \) is \( R \)-proper if for every condition \( p \in \forceP \), every \( \theta > 2^{|\forceP|} \), and every countable \( M \prec H(\theta) \) such that \( M \cap \omega_1 \in R \) and \( p, \forceP \in M \), there exists a condition \( q < p \) that is \( (M, \forceP) \)-generic. A condition \( q \in \forceP \) is said to be \( (M, \forceP) \)-generic if \( q \Vdash ``\dot{G} \cap M \) is an \( M \)-generic filter," where \( \dot{G} \) is the canonical name for the generic filter. See also \cite{Goldstern}.

\begin{lemma}
Let \( R \subset \omega_1 \) be stationary and co-stationary. Then the club-shooting forcing \( \forceP_R \) generically adds a club through \( R \). Moreover, \( \forceP_R \) is \( R \)-proper, \( \omega \)-distributive, and thus preserves \( \omega_1 \). Additionally, \( R \) and all its stationary subsets remain stationary in the generic extension.
\end{lemma}

We will select a family of sets \( R_{\beta} \) such that we can shoot an arbitrary pattern of clubs through the elements of \( R_{\beta} \), with this pattern being readable from the stationarity of the \( R_{\beta} \)'s in the generic extension. It is crucial to recall that for a stationary, co-stationary set \( R \subset \omega_1 \), \( R \)-proper posets can be iterated with countable support, and the result will again be an \( R \)-proper forcing. This follows from the well-known results for plain proper forcings (see \cite{Goldstern}, Theorem 3.9, and the subsequent discussion).

\begin{fact}
Let \( R \subset \omega_1 \) be stationary and co-stationary. Suppose \( (\forceP_{\alpha} : \alpha < \gamma) \) is a countable support iteration, and let \( \forceP_{\gamma} \) denote the resulting partial order obtained through the countable support limit. If at each stage \( \alpha \), \( \forceP_{\alpha} \Vdash \dot{\forceP}(\alpha) \) is \( R \)-proper, then \( \forceP_{\gamma} \) is \( R \)-proper.
\end{fact}

Once we decide to shoot a club through a stationary, co-stationary subset of \( \omega_1 \), this club will be contained in all \( \omega_1 \)-preserving outer models. This gives us a robust method of encoding arbitrary information into a suitably chosen sequence of sets, which has been used multiple times in previous work.

\begin{lemma}\label{coding with stationary sets}
Let \( (R_{\alpha} : \alpha < \omega_1) \) be a partition of \( \omega_1 \) into \( \aleph_1 \)-many stationary sets, let \( r \in 2^{\omega_1} \) be arbitrary, and let \( \forceP \) be a countable support iteration \( (\forceP_{\alpha} : \alpha < \omega_1) \), defined inductively by:
\[
\forceP(\alpha) := \dot{\forceP}_{\omega_1 \setminus R_{2 \cdot \alpha}} \text{ if } r(\alpha) = 1,
\]
and
\[
\forceP(\alpha) := \dot{\forceP}_{\omega_1 \setminus R_{(2 \cdot \alpha) + 1}} \text{ if } r(\alpha) = 0.
\]
Then in the resulting generic extension \( V[\forceP] \), we have for all \( \alpha < \omega_1 \):
\[
r(\alpha) = 1 \text{ if and only if } R_{2 \cdot \alpha} \text{ is nonstationary,}
\]
and
\[
r(\alpha) = 0 \text{ if and only if } R_{(2 \cdot \alpha) + 1} \text{ is nonstationary.}
\]
\end{lemma}

\begin{proof}
Assume without loss of generality that \( r(0) = 1 \). Then the iteration will be \( R_1 \)-proper, hence \( \omega_1 \)-preserving. Now let \( \alpha < \omega_1 \) be arbitrary, and assume that \( r(\alpha) = 1 \) in \( V[\forceP] \). By the definition of the iteration, we must have shot a club through the complement of \( R_{2 \cdot \alpha} \), so \( R_{2 \cdot \alpha} \) is nonstationary in \( V[\forceP] \).

Conversely, if \( R_{2 \cdot \alpha} \) is nonstationary in \( V[\forceP] \), assume for contradiction that we did not use \( \forceP_{\omega_1 \setminus R_{2 \cdot \alpha}} \) in the iteration \( \forceP \). Note that for \( \beta \neq 2 \cdot \alpha \), every forcing of the form \( \forceP_{\omega_1 \setminus R_{\beta}} \) is \( R_{2 \cdot \alpha} \)-proper, as \( \forceP_{\omega_1 \setminus R_{\beta}} \) is \( \omega_1 \setminus R_{\beta} \)-proper, and \( R_{2 \cdot \alpha} \subset \omega_1 \setminus R_{\beta} \). Hence, the iteration \( \forceP \) would be \( R_{2 \cdot \alpha} \)-proper, and the stationarity of \( R_{2 \cdot \alpha} \) would be preserved. But this leads to a contradiction.
\end{proof}


The second forcing technique we employ is the \emph{almost disjoint coding forcing}, introduced by R. Jensen and R. Solovay (see \cite{JensenSolovay}). In this context, we identify subsets of \( \omega \) with their characteristic functions, and we use the term \emph{reals} both for elements of \( 2^\omega \) and for subsets of \( \omega \), respectively.

Let \( D = \{ d_{\alpha} \mid \alpha < \aleph_1 \} \) be a family of almost disjoint subsets of \( \omega \), meaning that for any two sets \( r, s \in D \), the intersection \( r \cap s \) is finite. Let \( X \subset \kappa \) be a set of ordinals, where \( \kappa \leq 2^{\aleph_0} \). There exists a ccc forcing, denoted \( \mathbb{A}_D(X) \), which adds a new real \( x \) that encodes \( X \) relative to the family \( D \). The encoding satisfies the following condition:

\[
\alpha \in X \text{ if and only if } x \cap d_{\alpha} \text{ is finite.}
\]

\begin{definition}
The almost disjoint coding forcing \( \mathbb{A}_D(X) \), relative to an almost disjoint family \( D \), consists of conditions \( (r, R) \in \omega^{<\omega} \times D^{<\omega} \). The partial order \( (s, S) \le (r, R) \) holds if and only if:
\begin{enumerate}
  \item \( r \subseteq s \) and \( R \subseteq S \),
  \item If \( \alpha \in X \) and \( d_{\alpha} \in R \), then \( r \cap d_{\alpha} = s \cap d_{\alpha} \).
\end{enumerate}
\end{definition}

For the remainder of this paper, we let \( D \in L \) be the definable almost disjoint family of reals obtained by recursively adding the \( <_L \)-least real to the family, ensuring that each new real is almost disjoint from all previously chosen reals. Whenever we use almost disjoint coding forcing, we assume that we code relative to this fixed almost disjoint family \( D \).

The next two forcings we briefly discuss are Jech's forcing for adding a Suslin tree with countable conditions, and the associated forcing that adds a cofinal branch through a Suslin tree \( S \). Recall that a set-theoretic tree \( (S, <) \) is a \emph{Suslin tree} if it is a normal tree of height \( \omega_1 \) and has no uncountable antichain. Forcing with a Suslin tree \( S \), where conditions are just nodes in \( S \), is a ccc forcing of size \( \aleph_1 \). Jech’s forcing to generically add a Suslin tree is defined as follows:

\begin{definition}
Let \( \forceP_J \) be the forcing whose conditions are countable, normal trees ordered by end-extension, i.e. \( T_1 \le T_2 \) if and only if there exists \( \alpha \leq \operatorname{height}(T_1) \) such that \( T_2 = \{ t \upharpoonright \alpha \mid t \in T_1 \} \).
\end{definition}

It is well-known that \( \forceP_J \) is \( \sigma \)-closed and adds a Suslin tree. In fact, \( \forceP_J \) is forcing equivalent to \( \mathbb{C}(\omega_1) \), the forcing for adding a Cohen subset to \( \omega_1 \) with countable conditions. The tree \( T \) generically added by \( \forceP_J \) has the additional property that for any Suslin tree \( S \) in the ground model, \( S \times T \) will also be a Suslin tree in \( V[G] \). This property can be used to develop a robust coding method (see also \cite{Ho} for further applications).

\begin{lemma}\label{1 preservation of Suslin trees}
Let \( V \) be a universe, and let \( S \in V \) be a Suslin tree. Suppose \( \forceP_J \) is Jech’s forcing for adding a Suslin tree. Let \( g \subset \forceP_J \) be a generic filter, and let \( T = \bigcup g \) denote the generic tree. If we let \( T \in V[g] \) be the forcing that adds an \( \omega_1 \)-branch \( b \) through \( T \), then:

\[
V[g][b] \models S \text{ is Suslin.}
\]
\end{lemma}

\begin{proof}
Let \( \dot{T} \) be the \( \forceP_J \)-name for the generic Suslin tree. We claim that \( \forceP_J \ast \dot{T} \) has a dense subset which is \( \sigma \)-closed. Since \( \sigma \)-closed forcings preserve ground model Suslin trees, this will be sufficient. To prove the claim, consider the following set:

\[
\{ (p, \check{q}) \mid p \in \forceP_J \land \operatorname{height}(p) = \alpha+1 \land \check{q} \text{ is a node of } p \text{ at level } \alpha \}.
\]

It is easy to check that this set is dense and \( \sigma \)-closed in \( \forceP_J \ast \dot{T} \).
\end{proof}


The fact that the two-step iteration \( \forceP_J \ast \dot{T} \) contains a dense, \( \sigma \)-closed subset also implies that the following two-step iteration:

\[
\forceP_J \ast \prod_{i < \omega_1} \dot{T}
\]
which first adds a Suslin tree and then adds \( \omega_1 \)-many cofinal branches through the generic tree, using a countably supported product over \( V[\forceP_J] \), has a countably closed, dense subset. In fact, the set:

\begin{align*}
\{ (p, \vec{q}) \mid p \in \forceP_J, &\text{ height}(p) = \alpha + 1, \,  \\& \vec{q} \text{ is a countable sequence of nodes of } p \text{ at level } \alpha \}
\end{align*}
is dense in \( \forceP_J \ast \prod_{i < \omega_1} \dot{T} \) and \( \sigma \)-closed.



The Suslin trees added by Jech’s forcing have an important property that we will exploit. First, we define:

\begin{definition}
Let \( T \) be a Suslin tree. We say that \( T \) is \emph{Suslin off the generic branch} if, after forcing with \( T \) to add a generic branch \( b \), the tree \( T_p \) remains Suslin for every node \( p \in T \) that is not in \( b \).

More generally, let \( T \) be Suslin and let \( n \in \omega \). We say that \( T \) is \emph{n-fold Suslin off the generic branch} if, after forcing with \( T^n := \prod_{i < n} T \) (which adds \( n \)-many generic branches \( b_0, \dots, b_{n-1} \)), the tree \( T_p \) remains Suslin in \( V[T^n] \) for any node \( p \) that is not in any of the \( b_i \)'s.
\end{definition}

The next result, due to G. Fuchs and J. Hamkins, shows that Suslin trees added by Jech’s forcing have this property:

\begin{theorem}
Let \( T \) be a \( \forceP_J \)-generic Suslin tree. Then \( T \) is \emph{n-fold Suslin off the generic branch} for every \( n \in \omega \).
\end{theorem}

Note that if \( T \) is n-fold Suslin off the generic branch, and \( \{ t_i \mid i < \omega \} \) is an antichain in \( T \), then the countably supported iteration \( \Asterisk_{i < \omega} T_{t_i} \) is proper. Moreover, if \( t \) is a node that is not on any of the generic branches added by \( \Asterisk_{i < \omega} T_{t_i} \), then \( T_t \) remains Suslin in the generic extension by \( \Asterisk_{i < \omega} T_{t_i} \).

Similarly, if \( T \) is n-fold Suslin off the generic branch, then any \( \delta < \omega_1 \)-iteration of \( T \) with countable support is proper and thus preserves \( \aleph_1 \). Moreover, the full \( \omega_1 \)-length iteration \( \Asterisk_{i < \omega_1} T \) is proper.

This can be seen by noting that any condition \( \vec{p} \in (\Asterisk_{i < \delta} \check{T}) \) can be strengthened to a condition \( \vec{q} < \vec{p} \), such that there exists a set \( (\dot{t}_i \mid i < \delta) \), where each \( \dot{t}_i \) is a \( \Asterisk_{j < i} T \)-name for a node in \( T \), and:

\[
1 \Vdash_{\Asterisk_{i < \delta} T_i} \{ \dot{t}_i \mid i < \omega_1 \} \text{ is an antichain in } T.
\]

The forcing \( \Asterisk_{i < \delta} T \) below \( \vec{q} \) is equivalent to the forcing \( \Asterisk_{i < \delta} T_{\dot{t}_i} \). Since each \( T \) is n-fold Suslin off the generic branch, this forcing is a countable support iteration of forcings with the ccc, making it proper. Thus, \( \Asterisk_{i < \delta} \check{T} \) is locally proper, meaning that below any condition there is a stronger condition such that the forcing below it is proper. Local properness implies properness, as locally proper forcings preserve stationarity in $[\lambda]^{\omega}$ to which properness is equivalent.

W emphasize that we need countable support for this argument. A finite support iteration of \( \Asterisk_{i < \delta} \check{T} \) would collapse \( \aleph_1 \).



Next, we turn to Jech’s forcing \( \forceP_J \) and the product of \( \forceP_J \)-forcing. We observe that a product of \( \forceP_J \)-forcings will add Suslin trees, and we can destroy these trees without unwanted interference.

\begin{lemma}\label{ManySuslinTrees}
Let \( S \) be a Suslin tree in \( V \), and let \( \forceP \) be a countably supported product of length \( \omega_1 \) of forcings \( \forceP_J \), with \( G \) as its generic filter. Then in \( V[G] \), there is an \( \omega_1 \)-sequence of Suslin trees \( \vec{T} = (T_{\alpha} \mid \alpha \in \omega_1) \) such that for any finite \( e \subset \omega \) with pairwise different members, the tree \( S \times \prod_{i \in e} T_i \) will be a Suslin tree in \( V[G] \).
\end{lemma}



We define a sequence of Suslin trees as follows:

\begin{definition}
Let \( \vec{T} = (T_{\alpha} \mid \alpha < \kappa) \) be a sequence of Suslin trees. We say that the sequence is an \emph{independent family of Suslin trees} if, for every finite set of pairwise distinct indices \( e = \{ e_0, e_1, \dots, e_n \} \subset \kappa \), the product \( T_{e_0} \times T_{e_1} \times \cdots \times T_{e_n} \) is a Suslin tree again.
\end{definition}



We summarize the previous results as follows:

\begin{theorem}
Let \( \forceP \) be the countably supported product of Jech’s forcing \( \forceP_J \). Then \( \forceP \) adds an \( \omega_1 \)-sequence \( \vec{S} = (S_{\alpha} \mid \alpha < \omega_1) \) of independent Suslin trees with the following properties:
\begin{enumerate}
  \item Every \( S_{\alpha} \in \vec{S} \) is \( n \)-fold Suslin off the generic branch for every \( n \in \omega \),
  \item For every \( \alpha < \omega_1 \) and every \( S_{\alpha} \)-antichain of nodes \( \{ s_i \mid i < \delta < \omega_1 \} \subset S_{\alpha} \), the set:

\[
(\vec{S} \setminus \{ S_{\alpha} \}) \cup \{ (S_{\alpha})_{s_i} \mid i < \delta \}
\]

is an independent set of Suslin trees.
\end{enumerate}
\end{theorem}

\subsection{The ground model $W$ of the iteration}
We have to first create a suitable ground model $W$ over which the actual iteration will take place. $W$ will be a generic extension of $L$ which has no new reals. Moreover $W$ has the crucial property that in $W$ there is an $\omega_1$-sequence $\vec{S}$ of $\omega_1$ trees which is $\Sigma_1(\{\omega_1\})$-definable over $H(\omega_2)^W$ (i.e. the definiton is a $\Sigma_1$-formula with $\omega_1$ as its only parameter) and which forms an independent sequence of Suslin trees in an inner model of $W$. The sequence $\vec{S}$ will enable a coding method we will use throughout this article all the time.

To form $W$, we start with G\"odels constructible universe $L$ as our 
ground model.
We first fix an appropriate sequence of stationary, co-stationary subsets of $\omega_1$ as follows.
Recall that $\diamondsuit$ holds in $L$, i.e. over $L_{\omega_1}$ there is a 
$\Sigma_1$-definable sequence $(a_{\alpha} \, : \, \alpha < \omega_1)$ of countable subsets of $\omega_1$
such that any set $A \subset \omega_1$ is guessed stationarily often by the $a_{\alpha}$'s, i.e.
$\{ \alpha < \omega_1 \, : \, a_{\alpha}= A \cap \alpha \}$ is a stationary and co-stationary subset of $\omega_1$.
The $\diamondsuit$-sequence can be used to produce an easily definable sequence of stationary, co-stationary subsets: we list the reals in $L$ in an $\omega_1$ sequence $(r_{\alpha} \, : \, \alpha < \omega_1)$, and let $\tilde{r}_{\alpha} \subset \omega_1$ be the unique element of $2^{\omega_1}$ which copies $r_{\alpha}$ on its first $\omega$-entries followed by $\omega_1$-many 0's. Then, identifying $\tilde{r}_{\alpha} \in 2^{\omega_1}$ with the according subset of $\omega_1$, we define for every $\beta < \omega_1$
a stationary, co-stationary set in the following way:
\[R'_{\beta} = \{ \alpha < \omega_1 \, : \, a_{\alpha}= \tilde{r}_{\beta} \cap \alpha \}.\] It is clear that $\forall \alpha \ne \beta (R'_{\alpha} \cap R'_{\beta} \in \hbox{NS}_{\omega_1})$ and we obtain a sequence of pairwise disjoint stationary sets as usual via setting for every $\beta < \omega_1$ \[R_{\beta}= R'_{\beta} \backslash \bigcup_{\alpha < \beta} R'_{\alpha}.\] and let $\vec{R}=(R_{\alpha} \, : \, \alpha < \omega_1)$. Via picking out one element of $\vec{R}$ and re-indexing we assume without loss of generality that there is a stationary, co-stationary $R \subset \omega_1$, which has pairwise empty intersection with every $R_{\beta} \in \vec{R}$. 
Note that for any $\beta < \omega_1$, membership in $R_{\beta}$ is uniformly $\Sigma_1$-definable over the model $L_{\omega_1}$, i.e. there is a $\Sigma_1$-formula $\psi(x,y)$ such that for every $\beta < \omega_1$
$\alpha \in R_{\beta} \Leftrightarrow L_{\omega_1} \models \psi(\alpha, \beta)$.


We proceed with adding $\aleph_1$-many Suslin trees using of Jech's Forcing $ \forceP_J$. We let 
\[\forceQ^0 = \prod_{\beta \in \omega_1} \forceP_J \] using countable support. This is a $\sigma$-closed, hence proper notion of forcing. We denote the generic filter of $\forceQ^0$ with $\vec{S}=(S_{\alpha} \, : \, \alpha < \omega_1)$ and note that by Lemma \ref{ManySuslinTrees} $\vec{S}$ is independent.  We fix a definable bijection between $[\omega_1]^{\omega}$ and $\omega_1$ and identify the trees in $(S_{\alpha }\, : \, \alpha < \omega_1)$ with their images under this bijection, so the trees will always be subsets of $\omega_1$ from now on. 

We shall single out the first tree of $\vec{S}$, as later we will use this tree in a different way than all the other trees, namely to generically produce sets of indices where some coding will take place. For this reason we will re-index $\vec{S}= (S_{\alpha} \mid \alpha < \omega_1)$ in defining $S'_{-1}:= S_0$ and $\vec{S}':= (S_{\alpha} \mid 1 \le \alpha < \omega_1)$. To ease notation we will write $\vec{S}$ for the just defined $\vec{S}'$ again. That is \[\vec{S}= (S_{\alpha} \mid \alpha < \omega_1)\] is an independent sequence of Suslin trees and $S_{-1}$ is another Suslin tree such that
\[\vec{S} \cup S_{-1}\] still forms an independent sequence.


We work in $L[\forceQ^0]$ and will define the second block of forcings as follows: we add a generic branch through each element in $\vec{S}$, but add no branch through $S_{-1}$, that is we let
\[\forceQ^1= \prod_{\beta < \omega_1} S_{\beta}. \]
Note that by the argument from the proof of lemma \ref{1 preservation of Suslin trees}, this forcing has a dense subset which is $\sigma$-closed. Hence $L[\forceQ^0][\forceQ^1]$ is a proper and $\omega$-distributive generic extension of $L$.

In a third step we code the trees from $\vec{S} \cup S_{-1}$ into the sequence of $L$-stationary subsets $\vec{R}$ we produced earlier, using Lemma \ref{coding with stationary sets}. It is important to note, that the forcing we are about to define does preserve Suslin trees, a fact we will show later.
The forcing used in the third step will be denoted by $\mathbb{Q}^2$ and will itself be a countable support iteration of length $\omega_1 \cdot \omega_1$ whose components are countable support iteration themselves. First, fix a definable bijection $h \in L_{\omega_2}$ between $\omega_1 \times \omega_1$ and $\omega_1 \cup \{-1\}$ and write $\vec{R}$ from now on in ordertype $\omega_1 \cdot \omega_1$ making implicit use of $h$, so we assume that $\vec{R}= (R_{\alpha} \, : \, \alpha < \omega_1 \cdot \omega_1)$. We let $\alpha \in \omega_1 \cup \{ -1 \}$ and consider the tree $S_{\alpha} \subset \omega_1$. Defining the $\alpha$-th factor of our iteration $\forceQ^2$, we let $\forceQ^2 (\alpha)$ be the countable support iteration which codes the characteristic function of $S_{\alpha}$ into the $\alpha$-th $\omega_1$-block of the $R_{\beta}$'s just as in Lemma \ref{coding with stationary sets}. So $\forceQ^2 (\alpha)$ is a countable support iteration whose factors, denoted by $\forceQ^2(\alpha) (\gamma)$ are defined via
\[ \forall \gamma < \omega_1 \,(\forceQ (\alpha) (\gamma)= \dot{\forceP}_{\omega_1 \backslash R_{\omega_1 \cdot \alpha + 2 \gamma +1}}) \text{ if } S_{\alpha} (\gamma) =0 \]
and
\[ \forall \gamma < \omega_1 \, (\forceQ( \alpha) (\gamma)= \dot{\forceP}_{\omega_1 \backslash R_{\omega_1 \cdot \alpha + 2 \gamma}}) \text{ if } S_{\alpha} (\gamma) =1. \]


Recall that we let $R$ be a stationary, co-stationary subset of $\omega_1$ which is disjoint from all the $R_{\alpha}$'s which are used. It follows from Lemma \ref{coding with stationary sets} that for every $\alpha \in \omega_1 \cup \{-1\}$, $\forceQ^2 (\alpha)$ is an $R$-proper forcing which additionally is $\omega$-distributive.  Then we let $\mathbb{Q}^2$ be the countably supported iteration, $$\mathbb{Q}^2=\Asterisk_{\alpha \in  \omega_1 \cup \{-1\} } \forceQ^2 (\alpha)$$ which is again $R$-proper (and $\omega$-distributive as we shall see later).
This way we can turn the generically added sequence of trees $\vec{S}$ into a definable sequence of trees.
Indeed, if we work in $L[\vec{S}\ast \vec{b} \ast G]$, where $\vec{S} \ast\vec{b} \ast  G$ is $\forceQ^0 \ast \mathbb{Q}^1 \ast \forceQ^2$-generic over $L$, then, as seen in Lemma \ref{coding with stationary sets} 
\begin{align*}
\forall \alpha, \gamma < \omega_1 (&\gamma \in S_{\alpha} \Leftrightarrow R_{\omega_1 \cdot \alpha + 2 \cdot \gamma} \text{ is not stationary and} \\ &
\gamma \notin S_{\alpha} \Leftrightarrow  R_{\omega_1 \cdot \alpha + 2 \cdot \gamma +1} \text{ is not stationary})
\end{align*}
Note here that the above formula can be written in a $\Sigma_1(\{\omega_1\} )$-way (i.e. it can be written as a $\Sigma_1$ formula with the ordinal $\omega_1$ as the only parameter), as  it reflects down to $\aleph_1$-sized, transitive models of $\ZFP$ which contain a club through exactly one element of every pair $\{(R_{\alpha}, R_{\alpha+1}) \, : \, \alpha < \omega_1\}$.


Our goal is to use $\vec{S}$ for coding. For this it is essential, that the sequence remains independent in the inner  universe $L[\forceQ^0 \ast \forceQ^2]$. Note that this is reasonable as $\forceQ^0 \ast \forceQ^1 \ast \forceQ^2$ can be written as $\forceQ^0 \ast ( \forceQ^1 \times \forceQ^2)$, hence one can form the inner model $L[\forceQ^0 \ast \forceQ^2]$ without problems.

The following  line of reasoning is similar to \cite{Ho}.
Recall that for a forcing $\forceP$ and $M \prec H(\theta)$, a condition $q \in \forceP$ is $(M,\forceP)$-generic iff for every maximal antichain $A \subset \forceP$, $A \in M$, it is true that $ A \cap M$ is predense below $q$.
The key fact is the following (see \cite{Miyamoto2} for the case where $\forceP$ is proper, see e.g. \cite{Ho2} or \cite{Ho3} for a proof)
\begin{lemma}\label{preservation of Suslin trees}
 Let $T$ be a Suslin tree, $R \subset \omega_1$ stationary and $\forceP$ an $R$-proper
 poset. Let $\theta$ be a sufficiently large cardinal.
 Then the following are equivalent:
 \begin{enumerate}
  \item $\Vdash_{\forceP} T$ is Suslin
 
  \item if $M \prec H_{\theta}$ is countable, $\eta = M \cap \omega_1 \in R$, and $\forceP$ and $T$ are in $M$,
  further if $p \in \forceP \cap M$, then there is a condition $q<p$ such that 
  for every condition $t \in T_{\eta}$, 
  $(q,t)$ is $(M, \forceP \times T)$-generic.
 \end{enumerate}

\end{lemma}


In a similar way, one can show that Theorem 1.3 of \cite{Miyamoto2} holds true if we replace proper by $R$-proper for $R \subset \omega_1$ a stationary subset.
\begin{theorem}
Let $(\forceP_{\alpha})_{\alpha < \eta}$ be a countable support iteration of length $\eta$, let $R \subset \omega_1$ be stationary and suppose that for every $\alpha < \eta$, for the $\alpha$-th factor of the iteration $\dot{\forceP}(\alpha)$ it holds that $\Vdash_{\alpha}  ``\dot{\forceP}(\alpha)$ is $R$-proper and 
preserves every Suslin tree.$"$ Then $\forceP_{\eta}$ is $R$-proper and preserves every Suslin tree.
\end{theorem}
So in order to argue that our forcing $\forceQ^2$ preserves Suslin trees when used over the ground model $W[\forceQ^0]$, it is sufficient to show that every factor preserves Suslin trees.
This is indeed the case.
\begin{lemma}
Let $R \subset \omega_1$ be stationary, co-stationary, then the club shooting forcing $\forceP_R$ preserves Suslin trees.
\end{lemma}
We shall not prove the lemma and instead refer again to \cite{Ho1}, \cite{Ho2} or \cite{Ho3} for details.

Let us set $W:= L[\forceQ^0\ast \forceQ^1  \ast \forceQ^2 ]$ which will serve as our ground model for a second iteration of length $\omega_1$. To summarize the above:

\begin{theorem}\label{MainPropertiesOfW}
The universe $W= L[\forceQ^0 \ast \forceQ^1 \ast \forceQ^2]$ is an $\omega$-distributive generic extension of $L$, in particular no new reals are added and $\omega_1$ is preserved. Moreover $W$ satisfies the following list:
\begin{enumerate}
\item There  is a $\Sigma_1(\{\omega_1\})$-definable, independent sequence of  trees $\vec{S} \cup S_{-1}$ which are Suslin in the inner model $L[\forceQ^0][\forceQ^2]$, yet no tree, except $S_{-1}$ is Suslin in $W$.
\item The tree $S_{-1}$ is n-fold Suslin off the generic branch for every $n \in \omega$.
\item  For every antichain  $\{ s_i \in S \mid i < \delta \} \subset S_{-1}$, the set 
$\vec{S} \cup \{ (S_{-1} )_{s_i} \mid i < \delta \}$ is independent in $L[\forceQ^0][\forceQ^2].$
\end{enumerate}


\end{theorem}
\begin{proof}
To see that $W$ has no new reals uses a standard argument. As $\forceQ^0 \ast \forceQ^1$ does not add any reals it is sufficient to show that $\forceQ^2$ is $\omega$-distributive in $L[\forceQ^0][\forceQ^1]$.
Let $p \in \forceQ^2$ be a condition and assume that $p \Vdash ``\dot{r} $ is a countable sequence of ordinals$"$. We shall find a stronger $q < p$ and a set $r$ in the ground model such that $q \Vdash \check{r}=\dot{r}$. Let $M \prec H(\omega_3)$ be a countable elementary submodel which contains $p, \forceQ^2$ and $\dot{r}$ and such that $M \cap \omega_1 \in R$, where $R$ is our fixed stationary set from above. Inside $M$ we recursively construct a decreasing sequence $p_n$ of conditions in $\forceQ^2$, such that for every $n$ in $\omega,$ $p_n \in M$, $p_n$ decides $\dot{r}(n)$ and for every $\alpha$ in the support of $p_n$, the sequence $\operatorname{sup}_{n \in \omega} \operatorname{max}( p_n(\alpha))$ converges towards $M \cap \omega_1$ which is in $R$. Now, $q':= \bigcup_{n \in \omega} p_n$ and for every $\alpha< \omega_1$ such that $q'(\alpha)\ne 1$ (where 1 is the weakest condition of the forcing),  in other words for every $\alpha$ in the support of $q'$ we define $q(\alpha):= q'(\alpha) \cup \{(\omega,sup (M \cap \omega_1))\}$ and $q(\alpha)=1$ otherwise. Then $q=(q(\alpha))_{\alpha < \omega_1}$ is a condition in $\forceQ^2$, as can be readily verified and $q \Vdash \dot{r} = \check{r}$, as desired.

The first assertion should be clear from the above discussion. 
\end{proof}

%We note that, if we work over $W$ as our ground model and force with the countably supported product 
%$\prod_{i < \omega_1} S_{-1}$ in order to add $\aleph_1$-many cofinal branches to $S_{-1}$, then this will not add reals to $W$. Indeed, if we write $W=L[\forceQ^0][\forceQ^1][\forceQ^2]$ and consider $\prod_{i < \omega_1} S_{-1} \in W$, then we can re-arrange the resulting generic extension 
%\[ L[\forceQ^0][\forceQ^1][\forceQ^2] [\prod_{i < \omega_1} S_{-1}] \]
%and write it as
%\[ L[\forceQ^0] [\prod_{i < \omega_1} S_{-1}] [\forceQ^1] [\forceQ^2], \] and $\forceQ^0 \ast \prod_{i <\omega_1} S_{-1} \ast \forceQ^1$ has a dense, $\sigma$-closed subset so it does not add reals, and neither does $\forceQ^2$ when applied over $ L[\forceQ^0] [\prod_{i < \omega_1} S_{-1}] [\forceQ^1]$.


\begin{lemma}
Let $S$ be $n$-fold Suslin off the generic branch for every $n \in \omega$. 
Then any iteration of the form
\[ \forceP:= \Asterisk_{i < \delta < \omega_1} (\check{S} \ast \dot{\forceQ}_i) \]
where for every $i < \delta$, $\Vdash_{\forceP_i} ``\dot{\forceQ}_i$ is proper and preserves Suslin trees$"$, using countable support is a proper forcing.
\end{lemma}
\begin{proof}
The proof is almost the same to an argument we gave in the last section.
We use induction on the length $\delta$. If $\forceP= (\forceP(i) \mid i < \delta)$
is proper and we consider the next step $\check{S} \ast \dot{\forceQ}_i)$
of the iteration, then we first show that
$ \forceP \Vdash \check{S} \text{ is proper } .$
Indeed, if $G \subset \forceP$ is generic, and $s \in S$ is an arbitrary condition, then we can extend $s$ in $S$ to a $t$ such that $S_t$ is a Suslin tree in $V[G]$. So 
\[ \forceP \Vdash \check{S} \text{ is locally ccc} \]
where locally ccc just means that for any condition in the forcing there is an extension of such that forcing below that extension has the ccc. 

Any locally ccc forcing is proper, which can be seen easily if we consider the definition of proper using stationary set preservation (stationarity in $[\lambda]^{\omega}$).
So 
\[\forceP \Vdash \check{S} \text{ is proper }. \]

That $\dot{\forceQ}_i$ is proper follows from the assumption so the two step iteration is proper as claimed.


\end{proof}

We also note that proper forcings of a certain type will leave the sequence $\vec{S} = (S_{\alpha } \mid \alpha < \omega_1)$ independent.
\begin{lemma}\label{PullingOutATree}
Let $\vec{S} \ast \vec{b} \times G$ be an $L$-generic filter for $\forceQ^0 \ast \forceQ^1 \times \forceQ^2$.
Let $A \subset \omega_1$, $A \in L$, and let $\vec{b} \upharpoonright A = (b_{\alpha} \mid \alpha \in A)$.
Work in $W':=L[\vec{S} \ast (\vec{b} \upharpoonright A) \times G]$ and let $\forceP = (\forceP_{\alpha}, \forceP(\alpha) \mid \alpha < \delta < \omega_1) \in W' $ be 
a countable support iteration such that for every factor $\forceP(\alpha)$, either $\forceP(\alpha)$ is a $\forceP_{\alpha}$-name of $S_{-1}$, or $\forceP(\alpha)$ is a $\forceP_{\alpha}$-name of a proper, $\aleph_1$-sized forcing which preserves Suslin trees.

Let $\beta \notin A$. Then $S_{\beta} \in \vec{S}$ remains a Suslin tree in $W' [\forceP]$.
\end{lemma}
\begin{proof}
The forcing $\forceP$  preserves $\aleph_1$ by the last lemma.

We show the assertion via induction on the length of $\forceP$. If the length is 1, the lemma follows from the fact that $\vec{S} \cup S_{-1}$ is independent. 

Assume that the lemma is true for iterations of length $\delta$ and we want to show it for iterations $\forceP$ of lenght $\delta+1$. If $\forceP = \forceP_{\delta} \ast \forceP(\delta)$ and $\forceP(\delta)$ is the name of a Suslin tree preserving forcing then $S_{\beta}$ remains Suslin in $W' [\forceP_{\delta}]$ by the induction hypothesis and hence also in $W'[\forceP_{\delta} ] [\forceP(\delta)] = W' [\forceP]$, as desired.

So we assume that $ \forceP(\delta)^{G_{\delta}} =S_{-1}$. The generic branch $b$ through $S_{-1}$ will be different from all the generic branches we added with $\forceP_{\delta}$. Thus there is a node $s_{\delta+1}$ such that $s_{\delta+1} \in b$ and $s_{\delta+1}$ is not in any of the $\forceP_{\delta}$ added $S_{-1}$-branches. As $S_{-1}$ is n-fold Suslin off the generic branch for every $n$, forcing at stage $\delta+1$ with $S$ below the condition $s_{\delta+1}$ is a ccc forcing over $W[\forceP_{\delta}]$ and also ccc over $W' [\forceP_{\delta}]$ and, by theorem \ref{MainPropertiesOfW}, forcing with $S_{-1}$ below $s_{\delta+1}$  preserves ground model Suslin trees, so in particular it preserves $S_{\beta}$.



\end{proof}



We end with a straightforward lemma which is used later in coding arguments.

\begin{lemma}\label{a.d.coding preserves Suslin trees}
 Let $T$ be a Suslin tree and let $\mathbb{A}_D(X)$ be the almost disjoint coding which codes
 a subset $X$ of $\omega_1$ into a real with the help of an almost disjoint family
 of reals $D$ of size $\aleph_1$. Then $$\mathbb{A}_{D}(X) \Vdash_{} T \text{ is Suslin }$$
 holds.
\end{lemma}
\begin{proof}
 This is clear as $\mathbb{A}_{D}(X)$ has the Knaster property, thus the product $\mathbb{A}_{D}(X) \times T$ is ccc and $T$ must be Suslin in $V[{\mathbb{A}_{D}(X)}]$. 
\end{proof}






\section{Coding machinery}
We continue with the construction of the appropriate notions of forcing which we want to use in our proof. The goal is to first define a coding forcings $\operatorname{Code} (x)$ for reals $x$, which will force for $x$ that a certain $\Sigma^1_3$-formula $\Phi(x)$ becomes true in the resulting generic extension. The coding method is slightly different than the one in  \cite{Ho2}, \cite{Ho3} or \cite{Ho4} leading to a slicker presentation. Recall that when we talk about the $\vec{S}$-sequence we mean all the members with index $\ge 0$, and leave $S_{-1}$ out.

We will work over the model $W$. We shall define the coding forcing $\operatorname{Code} (x,y,m)$ now, which codes the triple $(x,y,m)$ into $\vec{S}$ and is defined as a two step iteration where the first factor is forcing with $S_{-1}$ and the second factor is the almost disjoint coding forcing of a specific set $Y \subset \omega_1$. Let $(x,y)$ be two reals in $W$ and let $m$ be a natural number. We simply write $(x,y,m)$ for a real $w$ which codes the triple $(x,y,m)$ in a recursive way.

In the first step we add an $\omega_1$-branch $b$ through $S_{-1}$ and work in $W[b]$. In a second step we will form a specific set $Y \subset \omega_1$ and let the second forcing be $\dot{\mathbb{A}} \dot{(Y)}.$
We will define the crucial set $Y \subset \omega_1$ now. We consider the branch $b$ through $S_{-1}$ as  $b \subset \omega_1$. 
 
  Let 
\begin{align*}
A:= &\{ \omega \gamma +2n \mid \gamma \in b, n \notin (x,y,m) \} \cup \\& \{\omega \gamma + 2n+1 \mid \gamma \in b, n \in (x,y,m) \}.
\end{align*}




Let $X \subset \omega_1$ be chosen such that it codes the following objects:
\begin{itemize}
\item The generically added branch $b \subset S_{-1}$
\item The set $A \subset \omega_1$.
\item For every ordinal $\omega \gamma +2n \in A$ or $\omega \gamma +2n +1 \in A$ we pick a generic branch $b_{\omega \gamma +2n}$ or $b_{\omega \gamma +2n+1}$ through the tree $S_{\omega \gamma +2n}$ or $S_{\omega \gamma +2n+1}$.
\item We also collect all $\aleph_1$-many, generically added clubs through $\vec{R}$ which are necessary to define the trees from $\vec{S} \cup S_{-1}$ using our $\Sigma_1 ( \{ \omega_1 \} )$-formula from the last section.
\end{itemize}

Note that, when working in $L[X]$ and if $\gamma \in b$ then
 we can read off $(x,y,m)$, and thus  we say that $(x,y,m)$ is coded into $\vec{S}$ at the $\omega$-block starting at $\gamma$,  via looking at the $\omega$-block of $\vec{S}$-trees starting at $\gamma$ and determine which tree has an $\omega_1$-branch in $L[X]$.
\begin{itemize}
 \item[$(\ast)_{\gamma}$]  $n \in (x,y,m)$ if and only if $S_{\omega \cdot \gamma +2n+1}$ has an $\omega_1$-branch, and $n \notin (x,y,m)$ if and only if $S_{\omega \cdot \gamma +2n}$ has an $\omega_1$-branch.
\end{itemize}
Indeed if $n \notin (x,y,m)$ then we added a cofinal branch through $S_{\omega \cdot \gamma+ 2n}$. If on the other hand $S_{\omega \cdot\gamma +2n}$ does not have an $\omega_1$-branch in $L[X]$ then we must have added an $\omega_1$-branch through $S_{\omega \cdot \gamma +2n+1}$ as we always add an $\omega_1$-branch through either $S_{\omega \cdot \gamma +2n+1}$ or $S_{\omega \cdot \gamma +2n}$ and adding branches through some $S_{\alpha}$'s  will not affect that some $S_{\beta}$ remain Suslin in $L[X]$, as $\vec{S}$ is independent.



We note that we can apply an argument resembling David's in this situation. We rewrite the information of $X \subset \omega_1$ as a subset $Y \subset \omega_1$ using the following line of reasoning.
It is clear that any transitive, $\aleph_1$-sized model $M$ of $\ZFP$ which contains $X$ will be able to first define $\vec{S}$ correctly  and also correctly decode out of $X$ all the information regarding $(x,y,m)$ being coded at each $\omega$-block of $\vec{S}$ starting at every $\gamma \in b$. 
Consequently, if we code the model $(M,\in)$ which contains $X$ as a set $X_M \subset \omega_1$, then for any uncountable $\beta$ such that $L_{\beta}[X_M] \models \ZFP$:
\[L_{\beta}[X_M] \models \text{\ldq The model decoded out of }X_M \text{ satisfies $(\ast)_{\gamma}$ for every $\gamma \in b$\rdq.} \]
In particular there will be an $\aleph_1$-sized ordinal $\beta$ as above and we can fix a club $C \subset \omega_1$ and a sequence $(M_{\alpha} \, : \, \alpha \in C)$ of countable elementary submodels  of $L_{\beta} [X_M]$ such that
\[\forall \alpha \in C (M_{\alpha} \prec L_{\beta}[X_M] \land M_{\alpha} \cap \omega_1 = \alpha)\]
Now let the set $Y\subset \omega_1$ code the pair $(C, X_M)$ such that the odd entries of $Y$ should code $X_M$ and if $E(Y)$ denotes  the set of even entries of $Y$ and $\{c_{\alpha} \, : \, \alpha < \omega_1\}$ is the enumeration of $C$ then
\begin{enumerate}
\item $E(Y) \cap \omega$ codes a well-ordering of type $c_0$.
\item $E(Y) \cap [\omega, c_0) = \emptyset$.
\item For all $\beta$, $E(Y) \cap [c_{\beta}, c_{\beta} + \omega)$ codes a well-ordering of type $c_{\beta+1}$.
\item For all $\beta$, $E(Y) \cap [c_{\beta}+\omega, c_{\beta+1})= \emptyset$.
\end{enumerate}
We obtain
\begin{itemize}
\item[$({\ast}{\ast})$] For any countable transitive model $M$ of  ``$\ZFP$ and $\aleph_1$ exists$"$ such that $\omega_1^M=(\omega_1^L)^M$ and $ Y \cap \omega_1^M \in M$, $M$ can construct its version of the universe $L[Y \cap \omega_1^N]$, and the latter will see that there is an $\aleph_1^M$-sized transitive model $N \in L[Y \cap \omega_1^N]$ which models $(\ast)$ for $(x,m,k)$ and every $\gamma \in b \cap M$.
\end{itemize}
Thus we have a local version of the property $(\ast)$.

We have finally defined the desired set $Y$ and now we use 
\(\mathbb{A}_D (Y) \)
 relative to our previously defined, almost disjoint family of reals $D \in  L $ (see the paragraph after Definition 2.5)  to code the set $Y$ into one real $r$. This forcing only depends on the subset of $\omega_1$ we code, thus $\mathbb{A}_D(Y)$ will be independent of the surrounding universe in which we define it, as long as it has the right $\omega_1$ and contains the set $Y$.

The effect of the coding forcing $\operatorname{Code} (x,y,m)$ is that it generically adds a real $r$ such that
\begin{itemize}
\item[$({\ast}{\ast}{\ast})$] For any countable, transitive model $M$ of ``$\ZFP$ and $\aleph_1$ exists$"$, such that $\omega_1^M=(\omega_1^L)^M$ and $ r  \in M$, $M$ can construct its version of $L[r]$, denoted by $L[r]^M$, which in turn thinks that there is a transitive $\ZFP$-model $N$ of size $\aleph_1^M$  such that $N$ believes $(\ast)$ for $(x,y,m)$ and every $\gamma \in b \cap M$.
\end{itemize}

Indeed, if $r$ and $M$ are as above, then $M$ and $L[r]^M$ will compute the almost disjoint family $D$ up to the real indexed with $\omega_1 \cap M$ correctly, as discussed below the definition 2.3. As a consequence, $L[r]^M$ will contain the set $Y \cap \omega_1^M$, where $Y \subset \omega_1$ is as in $(\ast\ast)$. So in $L[Y \cap \omega_1^M]$, there is an $\aleph_1^M$-sized, transitive $N$ which models $(\ast)_{\gamma}$ for every $\gamma \in h \cap M$, as claimed.


Note that $({\ast} {\ast} {\ast})$ is a $\Pi^1_2$-formula in the parameters $r$ and $(x,y,m)$, as the set $h \cap M \subset \omega_1^M$ is coded into $r$. We say in the above situation that the real $(x,y,m)$ \emph{ is written into $\vec{S}$}, or that $(x,m,k)$ \emph{is coded into} $\vec{S}$. To summarize our discussion, given an arbitrary real of the form $(x,y,m)$, then our forcing $\operatorname{Code} (x,y,m)$, when applied over $W$, will add a real $r$ which will turn the $\Pi^1_2$-formula $({\ast} {\ast}{\ast})$ for $r,(x,y,m)$ into a true statement in $W[\operatorname{Code} (x,y,m))]$.



 
 The projective and local statement $({\ast} { \ast} {\ast} )$, if true,  will determine how certain inner models of the surrounding universe will look like with respect to branches through $\vec{S}$.
That is to say, if we assume that $({\ast} { \ast} {\ast} )$ holds for a real $(x,y,m)$ and is the truth of it is witnessed by a real $r$. Then $r$ also witnesses the truth of $({\ast} { \ast} {\ast} )$ for any transitive  model $M$ of  the theory ``$\ZFP+$ $\aleph_1$ exists and $\aleph_1= \aleph_1^L"$,  which contains $r$ (i.e. we can drop the assumption on the countability of $M$).
Indeed if we assume 
that there would be an uncountable, transitive $M$, $r \in M$, which witnesses that $({\ast} { \ast} {\ast} )$ is false. Then by L\"owenheim-Skolem, there would be a countable $N\prec M$, $r\in N$ which we can transitively collapse to obtain the transitive $\bar{N}$. But $\bar{N}$ would witness that $({\ast} { \ast} {\ast} )$ is not true for every countable, transitive model, which is a contradiction.

Consequently, the real $r$ carries enough information that
the universe $L[r]$ will see that certain trees from $\vec{S}$ have branches in that
\begin{align*}
n \in (x,m,k) \Rightarrow L[r] \models  ``S_{\omega \gamma + 2n+1} \text{ has an $\omega_1$-branch}".
\end{align*}
and
\begin{align*}
n \notin (x,m,k) \Rightarrow L[r] \models ``S_{\omega \gamma + 2n} \text{ has an $\omega_1$-branch}".
\end{align*}
Indeed, the universe $L[r]$ will see that there is a transitive model $N$ of ``$\ZFP+$ $\aleph_1$ exists and $\aleph_1=\aleph_1^L"$ which believes $(\ast)$ for every $\gamma \in b\subset \omega_1$, the latter being coded into $r$. But by upwards $\Sigma_1$-absoluteness, and the fact that $N$ can compute $\vec{S}$ correctly, if $N$ thinks that some tree in $\vec{S}$ has a branch, then $L[r]$ must think so as well.

\section{Allowable forcings}

Next we define the set of forcings which we will use in our proof.
We aim to iterate the coding forcings we defined in the last section. The coding forcings should take care of two tasks simultaneously: on the one hand it should work towards $\Pi^1_3$-uniformization, on the other hand it should produce a $\Delta^1_3$-definable well-order of the reals. These two tasks are reflected in the way we define allowable, which uses two cases. 

 The second case should take care of producing a $\Delta^1_3$-definable well-order of the reals in the following way. If $\forceP$ is some iteration relative to some bookkeeping $F$, and if $F(\beta)$ is $(\dot{z}_0,\dot{z}_1)$ and 
$\dot{z}_0^{G_{\beta}}=z_0$ and $\dot{z}_1^{G_{\beta}}=z_1$, then we look at all the forcing names (not necessarily $\forceP_{\beta}$-names) $\sigma_0,\sigma_1$ such that 
\[z_0 =\sigma_0^{G_{\beta}} \land z_1 = \sigma_1^{G_{\beta}}.\]
We pick the $<$-least such names $\sigma_0$ and $\sigma_1$ in the canonical, global well-order of $L$.

Then we let $z_0 < z_1$ (we use $<$ here again and hope no confusion arises) if and only if the $<$-least such $\sigma_0$ is less than the $<$-least such $\sigma_1$.
Then allowable is just pre-allowable which codes up pairs of reals $(z_0,z_1)$ according to the just defined well-order $<$.

That is 

\begin{definition}
Let $W$ be our ground model. Let $\alpha < \omega_1$ and let $F\in W$, $F: \alpha \rightarrow W$ be a bookkeeping function.
A countable support iteration $\forceP=(\forceP_{\beta}\,:\, {\beta< \alpha})$ is called allowable (relative to the bookkeeping function $F$)  if the function $F: \alpha \rightarrow W$  determines $\forceP$ inductively as follows:
 \begin{itemize}
 \item We assume that $\beta \ge 0$ and $\forceP_{\beta}$ is defined.
 We let $G_{\beta}$ be a $\forceP_{\beta}$-generic filter over $W$ and assume that $F(\beta)=(\dot{x},\dot{y},\dot{m})$, for a tuple of $\forceP_{\beta}$-names. We assume that $\dot{x}^{G_{\beta}}=:x$ and $\dot{y}^{G_{\beta}}=:y $ are reals and $\dot{m}^{G_{\beta}}=m$ is a natural number.
 
 \begin{itemize}
 \item[]Then we let $\forceP(\beta):= \operatorname{Code} (x,y,m)= S_{-1} \ast \dot{\mathbb{A}} (\dot{Y} )$.
 \end{itemize}
 \item If on the other hand $F(\beta)= (\dot{z_0},\dot{z_1})$ for two $\forceP_{\beta}$-names of reals $\dot{z_0}$ and $\dot{z_1}$, and if we let $z_i = \dot{z}_i^{G_{\beta}}$ then we either force with 
 \begin{itemize}
 \item[] $\forceP(\beta):= \operatorname{Code} (z_0,z_1)= S_{-1} \ast \dot{\mathbb{A}} (\dot{Y} )$, or with 
 \item[] $\forceP(\beta):= \operatorname{Code} (z_1,z_0)= S_{-1} \ast \dot{\mathbb{A}} (\dot{Y} )$.
 \end{itemize}
 according to whether $\sigma_0 < \sigma_1$ or $\sigma_1 < \sigma_0$, where $\sigma_0$ and $\sigma_1$ are as in the discussion above.

 \end{itemize}
 
\end{definition}



As allowable forcings form the base set of an inductively defined shrinking process, they are also denoted by 0-allowable with respect to $F$ to emphasize this fact.
Informally speaking, the bookkeeping $F$ hands us at every step a real of the form $(x,y,m)$ or $(z_0,z_1)$, and uses a cofinal branch through $S_{-1}$ which gives rise to a subset of $b \subset \omega_1$ (these sets we will call \emph{coding areas}) which corresponds to the places where we code up the relevant branches through $\vec{S}$ to compute $(x,y,m)$ or $(z_0,z_1)$ or $(z_1,z_0)$ using the coding mechanism described in the previous section. The definition ensures that the coding areas are an almost disjoint family of $\omega_1$, i.e. each two distinct coding areas have countable intersection. The definition of allowable ensures that each such coding area is used at most once in an allowable forcing, which will imply that we will not accidentally code an unwanted reals into $\vec{S}$. 

Every allowable forcing $\forceP$ can be written over $W$ as $ (\Asterisk_{\beta<\delta} S_{-1} \ast \mathbb{A} (\dot{Y}_{\beta}))$.


As a remark we add that for any allowable forcing $\forceP \in W$ and any $G \subset \forceP$ generic, the forcing $\forceP^G$ can correctly  be defined already over a proper inner model of  $W[G]$. Or put differently, for any allowable $\forceP$ and any condition $p \in \forceP$, forcing with $\forceP$ below $p$ can be defined already over an inner model of $W$.  Indeed, as $\forceP$'s definition depends only on the names of reals listed by $F$ and on the branches we add through $S_{-1}$, which form an almost disjoint family in $\omega_1$, we see that $\forceP^G$ can be defined in$W[G]$ using a countable list of names of reals and additionally the $\aleph_1$-many branches through $L$-Suslin from $\vec{S}$ which are used to define the sets $Y \subset \omega_1$ which then get coded using $\mathbb{A}_D (Y)$.
So  there are always $\aleph_1$-many trees from $\vec{S}$ and $\aleph_1$-many coding areas which are not used when defining $\forceP^G$ over $W[G]$. In particular, trees from $\vec{S}$ which are not used in $\forceP^G$'s definition over $W[G]$ can be forced with after we forced with $\forceP^G$, by the fact that we can re-arrange factors of a product forcing. This will become useful when showing that we do not accidentally code up information when iterating allowable forcings.

The next notion will become useful later.

\begin{definition}
Let $\forceP$ be an allowable forcing over $W$ relative to $F: \delta_1 \rightarrow H(\omega_2)$, and let $\forceQ$ be an allowable forcing over $W$ relative to the bookkeeping $H: \delta_2 \rightarrow H(\omega_2)$. Then we say $\forceQ$ is an allowable extension of $\forceP$, denoted by $\forceQ \triangleright \forceP$, if $\delta_1 < \delta_2$ and $H \upharpoonright \delta_1 = F$.
\end{definition}



If $\forceP \in W$ is a forcing such that there is an $\alpha < \omega_1$ and an $F \in W$, $F: \alpha \rightarrow W$ such that $\forceP$ is allowable with respect to $F$, then we often just drop the $F$ and simply say that $\forceP \in W$ is allowable.

We prove next some properties of allowable forcings.
 \begin{lemma} 


\begin{enumerate}
\item If $\forceP=(\forceP(\beta) \, : \, \beta < \delta) \in W$ is allowable then for every $\beta < \delta$, $\forceP_{\beta} \Vdash| \forceP(\beta)|= \aleph_1$.
\item Every allowable forcing over $W$ is an $R$-proper forcing over $L$. Thus it preserves $\aleph_1$ and every countable set in an allowable extension can be covered by a countable set from $L$.
\item Every allowable forcing over $W$ preserves $\CH$. Furthermore, if $\forceP= (\forceP(\alpha) \, : \, \alpha < \omega_1) \in W$ is an $\omega_1$-length iteration such that each initial segment of the iteration is allowable over $W$, then $W[\forceP] \models \CH$.
\item Assume that $\forceP=(\forceP_{\beta} \mid \beta <\delta_0)$ is allowable relative to $F^0$ and $\forceQ=(\forceQ_{\beta} \mid \beta < \delta_1)$ is allowable relative to $F^1$. Then the product of $\forceP$ and $\forceQ$ densely embeds into $\forceP \ast \check{\forceQ}$ and this densely embeds into an allowable forcing relative to a bookkeeping $F$.
\end{enumerate}
\end{lemma}
\begin{proof}
The first, the second and the third assertion follow immediately from the definition modulo some well-known results, for the third item see Theorem 2.10 and theorem 2.12 in \cite{Abraham}.

The first assertion of the fourth item is well-known, and the second assertion follows from the fact that  every countable set of ordinals added by an $R$-proper forcing can be covered by a countable set of ordinals from the ground model.
 So in particular the two step iteration $\forceP \ast \check{\forceQ}$, can be densely embedded into a countable support iteration from the ground model.
 What is left is to find a bookkeeping $F$ which witnesses allowability.
 The bookkeeping is just the concatenation of $F^0$ and $F^1$, that is $F(\beta) = F^0 (\beta)$ for $\beta < \delta_0$ and $F(\beta+\delta_0)= F^1(\beta)$, and we identify  the check-$\forceP$-names for conditions in $\forceQ$ with conditions of $\forceQ$ in the usual way. The way we defined the well-order of the reals $z_0 < z_1$ or $z_1 <z_0$ using arbitrary names $\sigma_0$, $\sigma_1$ and the global well-order of $L$, ensures that if  we consider $z_0, z_1$ in $W[\forceP]$ and on the other hand in $W[\forceP \times \forceQ]$, both universes will agree on whether $z_0 < z_1$ or $z_1 < z_0$, thus $\forceP \times \forceQ$ is indeed allowable with respect to $F$.


\end{proof}

For the rest of this work, we keep to identify $\forceP \ast \check{\forceQ}$  with
$\forceP \times \forceQ$ and  $\forceP \ast \check{\forceQ}$-names with $\forceP \times \forceQ$ names without further mentioning it. 

Let $\forceP= (\forceP(\beta) \, : \, \beta < \delta)$ be an allowable forcing with respect to some $F \in W$.
The set of  (names of) reals which are enumerated by $F$ is dubbed the set of reals which are coded by $\forceP$. That is, for every $\beta$, if we let $\dot{x}_{\beta}$ be the (name) of a real  listed by $F(\beta)$ and if we let $G \subset \forceP$ be a generic filter over $W$ and finally if we let
$ \dot{x}_{\beta}^G =:x_{\beta}$,  then we say that
$\{ x_{\beta} \, : \, \beta < \alpha \}$ is the set of reals coded by $\forceP$ and $G$ (though we will suppress the $G$).

Next we show, that iterations of 0-allowable forcings will not add accidentally new elements to the set of reals defined by the $\Sigma^1_3$-formula $\Phi(x,y,m)$ which claims the existence of a real $r$ for which $({\ast} {\ast} {\ast})$ with parameters $r$ and the real $(x,y,m)$ is true:
\begin{align*}
 \Phi(x,y,m) \equiv \exists r  \forall M (&M \text{ is countable and transitive and } M \models \ZFP \\&\text{ and } \omega_1^M=(\omega_1^L)^M \text{ and }  r, (x,y,m) \in M  \rightarrow M \models \varphi (r,(x,y,m)) )
\end{align*}
where $\varphi(r,(x,y,m))$ asserts that in $M$'s version of $L[r]$, there is a transitive, $\aleph_1^M$-sized $\ZFP$-model which witnesses that $(x,y,m)$ is coded into $\vec{S}$.


\begin{lemma}\label{nounwantedcodes}
If $\forceP \in W$ is allowable, $\forceP=(\forceP_{\beta} \, : \, \beta < \delta)$, $G \subset \forceP$ is generic over $W$ and $\{ x_{\beta} \, : \, \beta < \delta\}$ is the set of reals which are coded by $\forceP$. Let $\Phi(v_0)$ be the distinguished formula from above. Then
in $W[G]$, the set of reals which satisfy $\Phi(v_0)$ is exactly 
$\{ x_{\beta} \, : \, \beta < \delta\}$.
\end{lemma}
\begin{proof}
Let $G$ be $\forceP$ generic over $W$, we work in $W[G]$. Let $g= (b_{\beta} \, : \, {\beta} < \delta)$ be the set of the $\delta$-many $\omega_1$-branches through $S_{-1}$ we use to determine the coding areas in  the factors of $\forceP$.  Note that the family $\{b_{\beta} \,: \, \beta < \delta \}$ forms an almost disjoint family of subsets of $\omega_1$, that is two distinct members of  $\{b_{\beta} \,: \, \beta < \delta \}$ always have countable intersection. Thus there is an $\eta < \omega_1$ such that for arbitrary distinct $\beta_1$, $\beta_2 < \delta$,  $\eta> b_{\beta_1}\cap b_{\beta_2}$.  


We assume for a contradiction, that there is a real $(x,y,m)$ which satisfies $\Phi ((x,y,m))$ but $(x,y,m) \notin \{x_{\beta} \mid \beta < \delta \}$.
So there is a real $r$ such that $( {\ast} {\ast} {\ast})$ is true for $r$ and $(x,y,m)$. By the discussion above, $r$ carries enough information so that it codes a coding area $b$ and for every $\gamma \in b$ the following holds true.
\begin{align*}
n \in (x,y,m) \Rightarrow L[r] \models  ``S_{\omega \gamma + 2n+1} \text{ has an $\omega_1$-branch}".
\end{align*}
and
\begin{align*}
n \notin (x,y,m) \Rightarrow L[r] \models ``S_{\omega \gamma + 2n} \text{ has an $\omega_1$-branch}".
\end{align*}


As $(x,y,m)$ is distinct from every $x_{\beta}$, there must be $\aleph_1$-many $\alpha > \eta$ and an $n \in \omega$ such that, without loss of generality,
\[L[r] \models ``S_{\omega \alpha + 2n+1} \text{ has an $\omega_1$-branch}",\]
yet for every $r_{\beta}$ such that $({\ast} {\ast} {\ast} )$ holds true for $r_{\beta} $ and $x_{\beta}$,
\[L[r_{\beta}] \models  ``S_{\omega \alpha + 2n+1} \text{ does not have an $\omega_1$-branch}".\]

We fix an $\alpha$ which is as above. We claim  that there is no real in $W[G]$ such that $W[G] \models L[r] \models ``S_{\omega \alpha + 2n+1}$ has an $\omega_1$-branch$"$, which will be the desired contradiction.


We show this by pulling the generic branch for the forcing $S_{\omega \alpha + 2n+1}$ out of the generic for the  forcing $\forceQ^0 \ast \forceQ^1 \times \forceQ^2 \ast \forceP$ over $L$ which produces $W[G]$. 
Recall that the $\forceP$-generic $G$ produces $\delta$-many cofinal branches $\{b_{\beta} \mid \beta < \delta\}$ through $S_{-1}$. We let $\vec{S} \cup \{S_{-1}\}
 \ast \vec{b} \ast H$ denote a $(\forceQ^0 \ast \forceQ^1 \ast \forceQ^2)$-generic filter over $L$. Working in $W[G]$ we consider 
 \begin{align*}
A:= &\{ \omega \gamma +2n \mid \exists \beta < \delta (\gamma \in b_{\beta} \land n \notin x_{\beta}) \} \cup \\& \{\omega \gamma + 2n+1 \mid \exists \beta< \delta (\gamma \in b_{\beta} \land n \in x_{\beta} )\}.
\end{align*}
Recall that any almost disjoint coding forcing $\mathbb{A} (Y)$ can be defined in an absolute way over any universe containing $Y\subset \omega_1$.
As $\omega\alpha +2n+1 \notin A$  we note that in
\[ W[G]= L[\vec{S} \cup \{ S_{-1} \} ] [\{ b_i \mid i < \omega_1] [H] [G]\]
the forcing $\forceP^G$, whose factors are just branches through $S_{-1}$ and some almost disjoint coding forcings determined by the sets of the $x_{\beta}$'s and $b_{\beta}$'s, can already be defined in the inner model
\[ L[\vec{S}] [ \{ b_i \mid i < \omega_1, i \ne \omega\alpha +2n+1 \} [H] [G] \]
as can be seen by induction on the length of the iteration $\forceP$, using the fact that
$S_{\omega \alpha+2n+1}$ is still Suslin in $L[\vec{S}]  [\{ b_i \mid i < \omega_1, i \ne \omega\alpha +2n+1 \} [H] [G] $.

As a consequence, we can exploit the properties of product forcings to rearrange and obtain the equality
\begin{align*}
L[\vec{S} ] [ \{ b_i \mid i < \omega_1 \} ]& [H] [G]= \\& L[\vec{S}]  [ \{ b_i \mid i < \omega_1, i \ne \omega\alpha +2n+1 \} ] [H] [G] [b_{\omega\alpha +2n+1}] 
\end{align*}




To ease notation we let $B= \{ i <\omega_1 \mid i \ne \omega \alpha +2n +1 \}$
We claim now that $S_{\omega\alpha +2n+1}$ is a Suslin tree in $L[\vec{S} ] [\{b_{\beta} \mid \beta \in B \} ] [H] [G]$. This is clear by lemma \ref{PullingOutATree}.


Note now that, as forcing with $S_{\alpha+2n+1}$ is $\omega$-distributive, the reals will not change when passing from $L[\vec{S} ] [b_{\beta} \mid \beta \in B] [H] [G]$ to
$L[\vec{S} ] [b_{\beta} \mid \beta \in B] [H] [G] [b_{\alpha}]=W[G]$. 
 But this implies that 
\[L[\vec{S} ] [b_{\beta} \mid \beta \in B] [H] [G] \models \lnot \exists r L[r] \models `` S_{\alpha+2n+1} \text{ has an $\omega_1$-branch}" \]
as the existence of an $\omega_1$-branch through $S_{\alpha+2n+1}$ in the inner model $L[r]$ would imply the existence of such a branch in $L[\vec{S} ] [b_{\beta} \mid \beta \in B] [H] [G]$. Further, as no new reals appear when passing to $W[G]$ we also get 
\[W[G] \models \lnot \exists r L[r] \models `` S_{\alpha+2n+1} \text{ has an $\omega_1$-branch}". \]
This is the desired contradiction and as $G$ was an arbitrary generic filter the lemma is proven.

\end{proof}

\section{$\alpha$-allowability}

We define the derivative acting on the set of allowable forcings over \( W \). Inductively, for an ordinal \( \alpha \) and any bookkeeping function \( F \in W \), we assume that the notion of \( \zeta \)-allowable with respect to \( F \) has already been defined for every \( \zeta < \alpha \). Specifically, this means that for each \( \zeta < \alpha \), we have already defined a set of rules that, in conjunction with a bookkeeping function \( F \in W \), produces the following over \( W \):

\begin{itemize}
\item An allowable forcing \( \forceP = \forceP_{\delta} = (\forceP_{\beta} \, : \, \beta < \delta) \in W \), which is the actual forcing used in the iteration. Let \( G_{\delta} \) denote a \( \forceP_{\delta} \)-generic filter over \( W \).
  
\item A set \( I = \dot{I}_{\delta}^{G_{\delta}} = \) \( \{ (\dot{x}^{G_{\delta}}, \dot{y}^{G_{\delta}}, \dot{m}^{G_{\delta}}, \dot{\gamma}^{G_{\delta}}) : \dot{m} , \) \( \dot{x}, \dot{y}, \dot{\gamma} \text{ are } \forceP \text{-names for elements of } \omega, 2^{\omega}, \omega_1 \} \). The set \( I \in W[G_{\delta}] \) contains potential values for the uniformizing function \( f \) that we want to define. Note that for a given \( x \) and \( m \), there may be several values \( (x, y_1, m, \xi_1), \dots, (x, y_n, m, \xi_n) \). We say that \( (x, y, m) \) has rank \( \xi \) if \( (x, y, m, \xi) \in I \). There can be multiple ranks for a given \( (x, y, m) \). The goal is to use the \( (x, y, m) \)'s with the minimal rank, and among those, choose the one with the least name according to the fixed well-order of $L$ in the background, ensuring the well-definedness of our choice.
\end{itemize}

Following our established terminology, if applying the rules for \( \eta \)-allowable forcings over \( W \) and \( F \in W \) results in the pair \( (\forceP, I) \in W \), we say that \( \forceP \) is \( \eta \)-allowable with respect to \( F \) (over \( W \)), or simply that \( \forceP \) is \( \eta \)-allowable if there exists an \( F \) and an \( I \) such that \( \forceP \) is \( \eta \)-allowable with respect to \( F \).


We now define the derivation of the \( <\alpha \)-allowable forcings over \( W \), which we call \( \alpha \)-allowable (again over \( W \)). The definition is a uniform extension of 0-allowability. 

A \( \delta < \omega_1 \)-length iteration \( \forceP = (\forceP_{\beta} : \beta < \delta) \in W \) is called \( \alpha \)-allowable over \( W \) (or relative to \( W \)) if it is recursively constructed using two ingredients. First, a bookkeeping function \( F \in W \), \( F : \delta \to W^3 \), where for each \( \beta < \delta \), we write \( F(\beta) = ((F(\beta)_0, F(\beta)_1, F(\beta)_2)) \) for the corresponding values of the coordinates. Second, a set of rules similar to those for 0-allowability, with two additional rules added at each step of the derivative, which determine, along with \( F \), how the iteration \( \forceP \) and the set of \( f \)-values \( I \) are constructed.

The infinite set of rules is defined as follows. Fix a bookkeeping function \( F \in W \), \( F : \delta \to W^3 \), for \( \delta < \omega_1 \). Assume we are at stage \( \beta \) of our construction and that, inductively, we have already created the following list of objects:

\begin{itemize}
\item The forcing \( \forceP_{\beta} \in W \) up to stage \( \beta \), along with a \( \forceP_{\beta} \)-generic filter \( G_{\beta} \) over \( W \).
  
\item The set \( I_{\beta} = \dot{I}_{\beta}^{G_{\beta}} = \{ (\dot{x}^{G_{\beta}}, \dot{y}^{G_{\beta}}, \dot{m}^{G_{\beta}}, \dot{\zeta}^{G_{\beta}}) : \dot{m}, \dot{x}, \dot{y}, \dot{\zeta} \text{ are } \forceP_{\beta} \text{-names for elements of } \omega, 2^{\omega}, \omega_1 \} \), containing potential values for the uniformizing function \( \dot{f}^{G_{\beta}}(m, \cdot) \). Initially, we set \( I_0 = \emptyset \).
\end{itemize}

The set of possible \( f \)-values will change as the iteration progresses. Specifically, values for \( f \) must be added when a new, lower-ranked value of \( \dot{f}^{G_{\beta}}(m, x) \) is encountered. 

Now, working in \( W[G_{\beta}] \), we define the next forcing \( \forceP(\beta) \) and possibly update the set of possible values for the uniformizing function \( f(m, x) \). Assume that \( F(\beta)_0 = (\dot{x}, \dot{y}, \dot{m}) \), and let \( A \subset \beta \), \( A \in W \), be such that \( \dot{x}, \dot{y}, \dot{m} \) are \( \forceP_A = \ast_{\eta \in A} \forceP(\eta) \)-names, where we require that \( \forceP_A \in W \) is a subforcing of \( \forceP_{\beta} \) (e.g., if \( A = \gamma < \beta \) and \( F(\beta)_0 \) lists \( \forceP_{\gamma} \)-names). Let \( G_A := G_{\beta} \upharpoonright A \). We then set \( x = \dot{x}^{G_A}, y = \dot{y}^{G_A} \), and \(m= \dot{m}^{G_A} \), and proceed as follows:

\begin{enumerate}
\item[(a)]  
  \begin{itemize}
  \item[] There exists an ordinal \( \zeta < \alpha+1 \), chosen to be minimal, such that:
  \item[] First, we collect all \( \forceP_A \)-names for reals \( \dot{a} \). For each \( \forceP_A \)-name \( \dot{a} \), we pick the \( < \)-least nice name \( \dot{b} \) such that \( \dot{a}^{G_{\beta}} = \dot{b}^{G_{\beta}} \), and collect these names \( \dot{b} \) into a set \( C \). We assume that there is a \( < \)-least nice \( \forceP_A \)-name \( \dot{y_0} \) in \( C \) such that \( \dot{y_0}^{G_A} = y_0 \), 
  \[
  W[G_{\beta}] \models (x, y_0) \in A_m
  \]
  and there is no \( \zeta \)-allowable forcing \( \forceR \triangleright \forceP_{\beta} \), \( \forceR \in W \), extending \( \forceP_{\beta} \) such that 
  \[
  W[G_{\beta}] \models \forceR / G_{\beta} \Vdash (x, y_0) \notin A_m.
  \]
  If this condition holds, we proceed as follows:
  \end{itemize}

  \begin{itemize}
  \item Assume that \( F(\beta)_2 = (\dot{x}, \dot{z}, \dot{m}) \) is a triple of \( \forceP_A \)-names, with \( \dot{x}^{G_A} = x \), \( \dot{z}^{G_A} = z \neq y_0 \), and \( \dot{m}^{G_A} = m \). We define:
  \[
  \forceP(\beta) := \text{Code}(x, z, m).
  \]
  If the bookkeeping function does not have the desired form, we choose the \( < \)-least names of the desired objects and use them to define the forcing. In this case, we pick the \( < \)-

least \( \forceP_A \)-name for a countable ordinal \( \dot{\eta} \), and let \( \dot{z} \) be the \( < \)-least \( \forceP_A \)-name of a real such that \( \dot{z}^{G_A} \neq y_0 \). Then:
  \[
  \forceP(\beta) := \text{Code}(x, z, m).
  \]
  We also set \( \forceP_{\beta+1} = \forceP_{\beta} \ast \forceP(\beta) \) and let \( G_{\beta+1} = G_{\beta} \ast G(\beta) \) be its generic filter.

  \item We assign a new value to \( f \), i.e., set \( f(m, x) := y_0 \) and assign the rank \( \zeta \) to the value \( (x, y_0, m) \) in \( W[G_{\beta+1}] \). We update \( I_{\beta+1}^{G_{\beta+1}} := I_{\beta}^{G_{\beta}} \cup \{ (x, y_0, m, \zeta) \} \).
  \end{itemize}

\item[(b)] If case (a) does not apply, i.e., for each \( \zeta < \alpha \) and each pair of reals, the pair can be forced out of \( A_m \) by a \( \zeta \)-allowable forcing extending the current one, we let the bookkeeping function \( F \) fully determine what to force. We assume that \( F(\beta)_1 \) is a \( \forceP_A \)-name for a countable ordinal \( \dot{\eta} \), and let \( \dot{\eta}^{G_A} = \eta \). We assume that \( F(\beta)_2 \) is a nice \( \forceP_A \)-name for a pair of reals \( (\dot{x'}, \dot{y_0}) \) such that \( \dot{x'}^{G_A} = x \). We define:
  \[
  \forceP(\beta) := \text{Code}(x, y_0, m).
  \]
  Let \( G(\beta) \) be a \( \forceP(\beta) \)-generic filter over \( W[G_{\beta}] \) and set \( G_{\beta+1} = G_{\beta} \ast G(\beta) \).

  We do not update the set \( I_{\beta} \) of preliminary values for \( f \), i.e., we set \( I_{\beta+1} := I_{\beta} \).

  Otherwise, we choose the \( < \)-least \( \forceP_A \)-names for the desired objects \( g_{\eta} \) and \( (x, z, m) \), and force with:
  \[
  \forceP(\beta) := \text{Code}(x, z, m).
  \]

\item[(c)] If \( F(\beta) = (\dot{z_0}, \dot{z_1}) \) for two \( \forceP_{\beta} \)-names of reals \( \dot{z_0} \) and \( \dot{z_1} \), and if we let \( z_i = \dot{z}_i^{G_{\beta}} \), we either force with:

  \begin{itemize}
  \item[] \( \forceP(\beta) := \text{Code}(z_0, z_1) = S_{-1} \ast \dot{\mathbb{A}}(\dot{Y}) \),
  \item[] or \( \forceP(\beta) := \text{Code}(z_1, z_0) = S_{-1} \ast \dot{\mathbb{A}}(\dot{Y}) \),
  \end{itemize}
  depending on whether \( \sigma_0 < \sigma_1 \) or \( \sigma_1 < \sigma_0 \), where \( \sigma_0 \) and \( \sigma_1 \) are as defined in the discussion of 0-allowability above.
\end{enumerate}

At limit stages \( \eta \) of \( \alpha + 1 \)-allowable forcings, we use countable support:
\[
\forceP_{\eta} := \text{inv} \, \lim (\forceP_{\nu} : \nu < \eta).
\]
Finally, we set:
\[
I_{\eta}^{G_{\eta}} := \{ (m, x, y, \zeta) : \exists \xi < \eta \, ((m, x, y, \zeta) \in I_{\xi}^{G_{\xi}}) \}.
\]
This concludes the definition of the rules for \( \alpha + 1 \)-allowability over the ground model \( W \). To summarize:


\begin{definition}
Assume that $F \in W$, $F: \eta \rightarrow W^3$ is a bookkeeping function and that $\forceP=(\forceP_{\beta} \, : \, \beta < \eta)$ and $I=(I_{\beta} \, : \, \beta < \eta)$ is the result of applying the above defined rules together with $F$ over $W$. Then we say that $(\forceP,I)$ is $\alpha+1$-allowable with respect to $F$ (over $W$). Often, $I$ is clear from context, and we will just say $\forceP$ is $\alpha+1$-allowable with respect to $F$. We also say that $\forceP$ is $\alpha+1$-allowable over $W$ if there is an $F$ such that $\forceP$ is $\alpha+1$-allowable with respect to $F$.
\end{definition}
We add a couple of remarks concerning the definition of $\alpha+1$-allowable:

\begin{itemize}

\item   Once a value \( f_m(x) \) of rank \( \zeta < \alpha \) is identified during an \( \alpha \)-allowable iteration, the pair \( (x, f_m(x)) \) will remain part of \( A_m \) in all subsequent outer models generated by a \( \zeta \)-allowable extension. This ensures that once a value is locked in for \( f_m(x) \), it is impossible for a further \( \zeta \)-allowable forcing to destroy the membership \( (x,f_m (x) ) \in A_m \).


\item In case (a) of the definition, we focus on finding a pair \( (x, y) \) that satisfies \( \varphi_m \) in the inner model \( W[G_A] \), and cannot be removed from \( A_m \) by any further \( < \alpha \)-allowable forcing. Such a pair is a candidate for \( f_m(x) \). However, there may be better candidates for \( f_m(x) \), which could appear later in the forcing iteration or in another inner model, where ``better$"$ means having a lower rank. The selection process first minimizes the rank \( \zeta \) of the potential value \( f(m,x) \), which makes it easier to show that the notion of \( \alpha \)-allowability strengthens as \( \alpha \) increases. After minimizing rank, the next step is to minimize the set of triples \( (\dot{x}, \dot{y}, \dot{m}) \) based on the \( < \)-least ordering.

\item  The definition of \( \alpha+1 \)-allowable imposes an additional restriction compared to \( \alpha \)-allowable. In case (a), it considers not only forcings that are \( \beta \)-allowable for \( \beta < \alpha \), but also those that are \( \alpha \)-allowable. This creates a natural progression in the definitional strength, and as \( \alpha \) increases, the set of \( \alpha \)-allowable forcings becomes smaller. This results in more pairs of reals \( (x, y) \in A_m \) that can no longer be eliminated by further \( \alpha \)-allowable forcings. Consequently, case (a) will apply more frequently, leading to more restrictions in the definition of \( \alpha \)-allowable as \( \alpha \) increases. 


\end{itemize}

\par \medskip


\begin{lemma}\label{shrinkinglemma}
Work in $W$. If $\forceP$ is $\beta$-allowable over $W$ and $\alpha < \beta$, then $\forceP$ is $\alpha$-allowable over $W$. Thus the sequence of $\alpha$-allowable forcings (over $W$) is decreasing with respect to the $\subset$-relation.
\end{lemma}
\begin{proof}
Let \( \alpha < \beta \), and suppose \( \forceP \) is a \( \beta \)-allowable forcing with a bookkeeping function \( F \in W \) that determines \( \forceP \) according to the rules outlined previously. We will demonstrate that there exists a bookkeeping function \( F' \in W \) such that \( \forceP \) can also be viewed as an \( \alpha \)-allowable forcing determined by \( F' \).

The key idea is to define \( F' \) so that it simulates the reasoning for a \( \beta \)-allowable forcing at every stage, even though it will ultimately define an \( \alpha \)-allowable forcing.



We start by setting \( F'(\eta) = F(\eta) \) for all stages \( \eta \) until we reach the first stage, denoted \( \gamma \), where the case applied under the \( \alpha \)-allowable rules differs from the case applied under the \( \beta \)-allowable rules. 

By the minimality of \( \gamma \), the iteration \( \forceP_{\gamma} \) must agree when considered as both \( \alpha \)-allowable and \( \beta \)-allowable at this stage. At \( \gamma \), applying the \( \beta \)-allowable rules results in case (a), while applying the \( \alpha \)-allowable rules results in case (b).

At this stage, we have \( F(\gamma)_0 = (\dot{x}, \dot{y}, m) \), and \( G_{\gamma} \) is the generic filter for \( \forceP \), so that \( x = \dot{x}^{G_{\gamma}} \) and \( y = \dot{y}^{G_{\gamma}} \). There is a potential \( f(m,x) \)-value with rank \( \le \beta \) in \( W[G_A] \), where \( A \subset \gamma \) is such that \( \dot{x} \) and \( \dot{y} \) are \( \forceP_A \)-names, and \( (x, y) \in W[G_A] \). However, there is no potential \( f(m,x) \)-value with rank \( \le \alpha \).

In this case, we have the following:
\begin{itemize}
\item[-] A quadruple \( (x, y_0, m, \xi) \in I_{\gamma+1} \), where \( \xi \in (\alpha, \beta] \), and \( (x, y_0) \in W[G_A] \), with \( A \subset \gamma \) such that \( \dot{x} \) and \( \dot{y}_0 \) are \( \forceP_A \)-names.
\item[-]  A forcing \( \forceP(\gamma) = \text{Code}(\dot{x}, z, m) \) for a real \( z \neq y_0 \).
\end{itemize}


We now define \( F'(\gamma) = (F'(\gamma)_0, F'(\gamma)_1, F'(\gamma)_2) \) such that \( F'(\gamma)_0 = F(\gamma)_0 \) and \( \forceP(\gamma) \) is guessed correctly by \( F'(\gamma) \). Specifically, we define:
\begin{itemize}
\item[-]  \( F'(\gamma)_1 \) as \( (x, z, m) \),
\item[-] \( F'(\gamma)_2^{G_{\gamma}} := \eta \).
\end{itemize}


Note that the definition of \( F' \) is entirely within \( W \), and the use of \( G_{\gamma} \) is uniform, so it can be eliminated in the usual way.



The key observation is that when applying the rules for \( \alpha \)-allowable forcings at stage \( \gamma \) using \( F' \), the results will match those from the \( \beta \)-allowable iteration using \( F \). 

Thus, for the least stage \( \gamma \) where different cases apply when using \( F \) for \( \alpha \)-allowable and \( \beta \)-allowable forcings, we can define an \( F' \) such that \( \forceP_{\gamma+1} \) is an \( \alpha \)-allowable iteration, while it is also a \( \beta \)-allowable iteration using \( F \). 

This argument can be iterated. After handling \( (m, x, y) \) at stage \( \gamma \), we move to the next stage \( \gamma' > \gamma \) where the rules for \( \beta \)-allowable forcings (using \( F \)) yield a different case than those for \( \alpha \)-allowable forcings (using \( F' \)). We apply the same reasoning to show that we can continue treating the iteration as \( \alpha \)-allowable as we proceed with \( \forceP \), handling new triples in the same way. Thus, there exists a bookkeeping function \( F' \) such that \( \forceP \) is \( \alpha \)-allowable with respect to \( F' \).


\end{proof}

\begin{lemma}\label{alphaallowableclosedunderproducts}
Let $F_1 : \delta_1 \rightarrow W^3$ be a bookkeeping function which determines an $\alpha$-allowable forcing $\forceP^1=(\forceP^1_{\beta} \, : \,  \beta < \delta_1)$. Likewise let $F_2 : \delta_2 \rightarrow W^3$ be a bookkeeping function which determines an $\alpha$-allowable forcing $\forceP^2=(\forceP^2_{\beta} \, : \,  \beta < \delta_2)$.
Then the product $\forceP^1 \times \forceP^2$ is $\alpha$-allowable relative to a bookkeeping function $F \in W$ which is definable from $F_1$ and $F_2$.
\end{lemma}
\begin{proof}
We will prove it by induction on $\alpha$. For $\alpha=0$ the lemma is true.

Now assume that the lemma is true for $\alpha$. We shall argue that it is also true for $\alpha+1$.

We shall define a bookkeeping function $F'$ such that $\forceP^1 \times \forceP^2$ is $\alpha+1$-allowable relative to $F'$.
For ordinals $\beta < \delta_1$ we let $F'(\beta)=F_1(\beta)$. Then the $(\alpha+1)$-allowable forcing which will be produced on the first $\delta_1$-many stages is $\forceP^1$.

For $\beta> \delta_1$, we let $F'(\beta)= F_2( \beta -\delta_1)$. Then we claim that $F' \upharpoonright [\delta_1,\delta_1+\delta_2)$ using the rules of $(\alpha+1)$-allowability will produce $\forceP^2$.

First we prove by induction on $\beta \in [\delta_1, \delta_1 +\delta_2)$ that if $\beta$ is a stage such that case (b) applies when building the forcing using $F'$ over $W$, then case (b) also must apply at $\beta - \delta_1$ when building $\forceP^2$ over $W$ using $F_2$ and vice versa.

Assume first that at stage ${\beta-\delta_1}$, $\forceP^2_{\beta-\delta_1}$ is defined and case (b) applies there. That is,
if $F( {\beta-\delta_1})_0= (\dot{x}, \dot{y}, \dot{m})$ and  $G^2_{\beta-\delta_1}$ is our $\forceP^2_{\beta-\delta_1}$-generic filter over $W$, for every $\zeta < \alpha+1$ there is a $\zeta$-allowable $\forceQ  \triangleright \forceP^2_{\beta-\delta_1}$ such that \[ \forceQ \slash G^2_{\beta-\delta_1} \Vdash (\dot{x}^{ G^2_{\beta-\delta_1}},\dot{y}^{ G^2_{\beta-\delta_1}} ) \notin A_{\dot{m}^{G^2_{\beta-\delta_1}}}. \]
But then also $\forceP^1 \times \forceQ \triangleright \forceP^1 \times \forceP^2_{\beta}$ is $\zeta$-allowable, by the induction hypothesis and the previous lemma and 
\[ \forceP^1 \times (\forceQ \slash G^2_{\beta}) \Vdash (\dot{x}^{G^2_{\beta}} ,\dot{y}^{G^2_{\beta}}) \notin A_{\dot{m}^{G^2_{\beta}}} \] by upwards absoluteness of $\Sigma^1_3$-formulas. Thus we must be in  case (b) at stage $\beta$ as well, when defining $\forceP^1 \times \forceP^2$ using $F'$.

On the other hand if we are at stage $\beta \in [\delta_1, \delta_1+\delta_2)$, $F({\beta})_0= (\dot{x},\dot{y}, \dot{m})$, $G^1 \times G^2_{\beta-\delta_1}$ is an arbitrary $\forceP^1 \times \forceP^2_{\beta-\delta_1}$-generic filter over $W$ and we are in case (b) when defining the next forcing using $F'$, then, for every $\zeta < \alpha+1$ there is a $\zeta$-allowable $\forceR^{\zeta}  \triangleright \forceP^1 \times \forceP^2_{{\beta-\delta_1}}$ such that
\[ \forceR^{\zeta}\slash (G^1 \times G^2_{\beta-\delta_1} )\Vdash ({x},{y}) \notin A_m.\]  But then $\forceR^{\zeta} \triangleright \forceP^2_{{\beta-\delta_1}}$ is true, and the set  of$\forceR^{\zeta}$'s witnesses that we must be in case (b) at stage $\beta-\delta_1$ as well, when defining $\forceP^2$ using $F^2$ working over $W[G^2_{\beta-\delta_1}]$.


If we are at case (c) when defining our  $\alpha$-allowable forcing using $F'$ at stage $\beta$, then we must also be in case (c) when building $\forceP^2$ using $F^2$, as the definition just refers to a definable well-order of names, defined in $L$ and thus does not depend on the actual  model where we consider case (c). This reasoning also applies for the other direction, when defining an $\alpha$-allowable forcing over $W$ using $F^2$ and being in case (c), then we force with the very same forcing when
defining the $\alpha$-allowable forcing at stage $\delta_1 +\beta$ using $F'$. 


As a consequence we must be in the same cases when defining $\forceP^1 \times \forceP^2$ at stage $\beta$ over $W$ and when defining $\forceP^2$ at stage $\beta -\delta_1$ using $F_2$ over $W$. But then we let $F'(\beta)$ be such that it does exactly what $F_2(\beta-\delta_1)$ does. This implies that $\forceP^1 \times \forceP^2_{\beta+1}$ is $(\alpha+1)$-allowable with respect to $F' \upharpoonright \beta+1$ and the induction step is proven.

For $\beta$ being limit there is noting to show as the $\beta$-th forcing is uniquely determined by $\forceP_{\beta'} , \beta' < \beta$. 
Thus $F'$ witnesses that $\forceP^1 \times \forceP^2$ is $\alpha+1$-allowable.

The case where $\alpha$ is limit is identical to the induction step $\alpha \rightarrow \alpha+1$, modulo the obvious notational changes.


\end{proof}

It follows that once we encounter a potential $f_m(x)$-value $y$ of rank $<\alpha$, in an $\alpha$-allowable iteration, that $(x,y)$ will remain in $A_m$ for the rest of the $\alpha$-allowable iteration and for all further $\alpha$-allowable extensions: 
\begin{lemma}\label{fvaluesremain2}
Let $(\forceP,I)$ be $\alpha$-allowable over $W$ with respect to $F$ of length $\eta < \omega_1$, let $G$ be $\forceP$-generic and let $(x,y,m,\xi) \in I$ for some $\xi < \alpha$. Then in $W[G]$, $(x,y) \in A_m$ and for every $\forceQ \triangleright \forceP$ such that $\forceQ$ is $\xi$ allowable it holds that $\forceQ \slash G \Vdash (x,y) \in A_m$.
\end{lemma}
\begin{proof}
Let $\beta$ be the least stage in $\forceP$ such that $(x,y,m,\xi)$ is added to $I_{\beta}$. Then, as $\xi < \alpha$, we must be in case (a) at stage $\beta$.
This means that we found a pair $(x,y)$ which can not be kicked out of $A_m$ with a further $\xi$-allowable forcing, for a $\xi < \alpha$ as in the lemma. The tail however is an $\alpha$-allowable forcing over $W[G_{\beta}]$, hence also $\xi$-allowable and thus $(x,y) \in A_m$ throughout the tail of the iteration.

Adding an additional $\forceQ$ which is $\xi$-allowable to the tail of $\forceP$ does not alter the argument, which proves the second assertion of the lemma.
\end{proof}

\begin{lemma}
Work in $W$. For any $\alpha$, the set of $\alpha$-allowable forcings is nonempty.
\end{lemma}
\begin{proof}
By induction on $\alpha$. If there are $\alpha$-allowable forcings over $W$, then every bookkeeping function $F \in W$, $F: \delta \rightarrow W^3$  together with the rules (a) and (b) will create a nontrivial $\alpha+1$-allowable forcing just by the way we chose to define $\alpha+1$-allowability. For limit ordinals, the same reasoning applies.
\end{proof}

As a direct consequence we obtain that there must be an ordinal $\alpha$ such that for every $\beta> \alpha$, the set of $\alpha$-allowable forcings over $W$  must equal the set of $\beta$-allowable forcings over $W$. Indeed every allowable forcing is an $\aleph_1$-sized partial order in $W$, thus there are only set-many of them (modulo isomorphism), and the decreasing sequence of $\alpha$-allowable forcings must eventually stabilize at a set which also must be non-empty.

\begin{definition}
Let $\alpha_0$ be the least ordinal such that for every $\beta> \alpha_0$, the set of $\alpha_0$-allowable forcings over $W$ is equal to the set of $\beta$-allowable forcings over $W$. 
\end{definition}
The set of $\infty$-allowable forcings can also be described in the following way. A $\delta < \omega_1$-length iteration $\forceP= (\forceP_{\alpha} \, : \, \alpha< \delta)$ is $\infty$-allowable if it is recursively constructed following a bookkeeping function $F$ and a modified version of the two rules from above: we ask in (a) whether there exists an ordinal $\zeta$ at all for which the antecedens of (a) is true. If there is such an ordinal $\zeta$ we proceed as described in (a) if not we use (b). 
Note that for  $m\in \omega$ and a real $x$ we will have several potential $y$'s such that $(x,y,m,\xi) \in I$ as we go along in an $\infty$-allowable iteration. The ranks of the potential values form a decreasing sequence of ordinals, thus, once we set a value $f(m,x)$, we can be sure that eventually there will be a value for $f(m,x)$ which will not change any more in rank. 

\section{Definition of the universe in which the ${\Pi^1_3}$ uniformization property holds}
The notion of $\infty$-allowable will be used now to define the universe in which the ${\Pi^1_3}$-uniformization property is true and whose reals admit a $\Delta^1_3$-definable well-order. We let $W$ be our ground model and start an $\omega_1$-length iteration whose initial segments are all $\infty$-allowable with respect to $W$. We are using the following rules in combination with some bookkeeping $F \in W$. 
The actual properties of $F$ are not really relevant, $F$ should however satisfy that \begin{itemize}
\item $F: \omega_1 \rightarrow H(\omega_1)$ is surjective.
\item For every $x \in H(\omega_1)$, the set $F^{-1} (x)$ should be unbounded in $\omega_1$.
\end{itemize}


Inductively we assume that we are at stage $\beta < \omega_1$ of our iteration and the allowable forcings $\forceP_{\beta}$, $\forceR_{\beta}$ have been defined already. We assume additionally that the value $F(\beta)=(F(\beta)_0, F(\beta)_1)$ is in fact a pair of elements in $H(\omega_1)$ and $F(\beta)_0=(\eta_1, \eta_2, m)$ where $\eta_1 \le \beta$ and $\eta_2$ are ordinals and $m \in \omega$.
We let $(\dot{x}, \dot{y})$ be the $\eta_2$-th nice $\forceP_{\eta_1}$ name of a pair of reals relative to our wellorder $<$.
If we set  $\dot{x}^{G_{\beta}}=x$, $\dot{y}^{G_{\beta}}=y$ then we further assume that $W[G_{\beta}] \models (x,y) \in A_m$. Recall that $\alpha_0$ is the least ordinal such that the notion of $\alpha$-allowable stabilizes. We split into cases, following the definition of $\alpha_0$-allowable.

\begin{enumerate}
\item[(1)] We work in $W[G_{\eta_1}]$ and assume the following.
\begin{itemize}
\item There is an ordinal $\zeta \le \alpha_0$, which is chosen to be minimal for which
\item there is a $<$-least pair of nice $\forceP_{\eta_1}$-names $(\dot{x}',\dot{y}')$ such that $\dot{x}'^{G_{\eta_1}}=x$ and $\dot{y}'^{G_{\eta_1}}=y_0$ and $W[G_{\beta}] \models (x,y_0) \in A_m$ and
for every $\zeta$-allowable $\forceR \in W[G_{\beta}]$, 
\[ \forceR \Vdash (x,y_0) \in A_m. \]

\end{itemize}
If this is the case, then we set the following:
\begin{itemize}
\item We pick the $<$-least triple of $\forceP_{\eta_1}$-names $(\dot{x}, \dot{z}, m)$ such that if $\dot{x}^{G_{\eta_1}} =x$, $\dot{z}^{G_{\eta_1}} =z \ne y_0$ and such that $(x,z,m)$ has not been coded yet into $\vec{S}$ by $\forceP_{\beta}$. We let $\forceP(\beta):=\operatorname{Code} {(x,z,m)}$.  We also let $\forceP_{\beta+1} = \forceP_{\beta} \ast \forceP(\beta)$ and let $G_{\beta+1}$ be its generic filter.

\item  We set a new $f$ value, i.e. we set $f(m,x):=y'$ and assign the rank $\xi$ to the value. We update $I_{\beta+1}^{G'_{\beta+1}}:= I_{\beta}^{G'_{\beta}} \cup \{ (x,y',m,\xi)\}$.

\end{itemize}

\item [(2)]
We assume that case (a) is not true. 
So for every $\xi \le \alpha_0$, in particular for $\xi=\alpha_0$, every pair $(x,z) \in W[G_{\eta_1}]$ (so in particular for the pair $(x,y)$) 
there is a further $\alpha_0$-allowable $(\forceP^{(x,y)} \triangleright \forceP_{\beta}$, such that 
\begin{align*}
\forceP^{x,y} \slash G_{\beta} \Vdash (x,y) \notin A_m
\end{align*}

We pick the $<$-least such $\alpha_0$-allowable forcing $\forceP^{(x,y)}$ and use the tail $\forceP^{x,y} \slash G_{\beta}$ at stage $\beta$, and a generic filter $G^{x,y} \triangleright G_{\beta}$ for $\forceP^{x,y} \slash G_{\beta}$ over $W$ and set $G_{\beta+1}:= G^{x,y}$.

Then we update $I_{\beta+1}$ to be $I_{\beta} \cup I^{x,y}$. 

\item[(3)] If on the other hand $F(\beta)= (\dot{\eta_1},\dot{\eta_2})$ for two $\forceP_{\beta}$-names of countable ordinals $\dot{\eta_1}$ and $\dot{\eta_2}$ $\eta_i=\dot{\eta}_i^{G_{\beta}}$, and if we assume that $(z_0,z_1)$ are the 
$\eta_2$-th of a pair of  reals in $W[G_{\eta_1}]$, where $G_{\eta_1}$ is $G_{\beta}$ restricted to $\forceP_{\eta_1}$, then we let $z_i = \dot{z}_i^{G_{\beta}}$ and we either force with 
 \begin{itemize}
 \item[] $\forceP(\beta):= \operatorname{Code} (z_0,z_1)= S_{-1} \ast \dot{\mathbb{A}} (\dot{Y} )$, or with 
 \item[] $\forceP(\beta):= \operatorname{Code} (z_1,z_0)= S_{-1} \ast \dot{\mathbb{A}} (\dot{Y} )$.
 \end{itemize}
 according to whether $\sigma_0 < \sigma_1$ or $\sigma_1 < \sigma_0$, where $\sigma_0$ and $\sigma_1$ are as in the discussion above.


 
\end{enumerate}
As always we use countable support. This ends the definition of our iteration $((\forceP_{\beta}, I_{\beta}) \, : \, \beta < \omega_1)$. We set $\forceP_{\omega_1}$ to be the direct limit of $(\forceP_{\beta} \, : \, \beta < \omega_1)$, and $I_{\omega_1}=\bigcup_{\beta < \omega_1} I_{\beta}$.
\par \medskip
We next derive its basic properties. First we note that the iteration is such that there is an $F' \in W$, $F': \omega_1 \rightarrow W$, and such that for every $\delta < \omega_1$,  $((\forceP_{\beta}, I_{\beta}) \, : \, \beta < \delta)$ is $\alpha_0$-allowable over $W$ with respect to $F' \upharpoonright \delta$. Indeed the bookkeeping $F$ can be used to readily derive such an $F' \in W.$

\begin{fact}
The just defined iteration $(\forceP_{\omega_1},I_{\omega_1}) \in W$ is such that every initial segment is $\alpha_0$-allowable over $W$ relative to a fixed $F' \in W$.
\end{fact}
As a consequence, the $f_m$-values of rank $< \alpha_0$ we define as we go along the iteration are such that they will certainly belong to $A_m$ in the final model by Lemma \ref{fvaluesremain2}. We let $G_{\omega_1}$ be $\forceP_{\omega_1}$-generic over $W$. What is left, is to show that in $W[G_{\omega_1}]$, for every $m \in \omega$ and every real $x$ such that $A_{m,x} \ne \emptyset$, we do have exactly one pair of reals $(x,y) \in A_m $ such that $(x,y,m)$ is not coded into $\vec{S}$.
The next lemma does exactly that, and is the main step in proving that the $\Pi^1_3$-uniformization property holds true in $W[G_{\omega_1}]$.

\begin{lemma}
In $W[G_{\omega_1}]$ the following dichotomy holds true:
\begin{enumerate}
\item Either $(x,m)$ is such that there is a real $y$ and $\xi <  \alpha_0$ such that $(x,y,m,\xi)\in I$. Then there is a unique real $y_0$ such that
\[W[G_{\omega_1}] \models`` (x,y_0) \in A_m \land (x,y_0,m) \text{ is not coded somewhere into } \vec{S}". \]

\item Or $(x,m)$ is such that for every real $y$, if $\xi< \alpha_0$ then $(x,y,m,\xi) \notin I$, in which case 
\[W[G_{\omega_1}] \models `` \text{The $x$-section of $A_m$ is empty}" \]
\end{enumerate}
\end{lemma}
\begin{proof}
We assume first that the assumptions of case 1 are true, i.e. there is a $y$ and $\xi < \alpha_0$ such that $(x,y,m,\xi) \in I$. Then there is a real $y_0 \in W[G_{\omega_1}]$ (and an attached ordinal $\xi_0 < \alpha_0$) whose
$\forceP_{\omega_1}$-name is $<$-minimal among all such names.
We let $\beta$ be the least stage where we add $(x,y_0,m,\xi_0)$ to $I_{\beta}$.
\par \medskip

\begin{claim}
$W[G_{\omega_1}] \models (x,y_0) \in A_m$.
\end{claim}
\begin{proof}[Proof of the first Claim]
This follows immediately from the lemma \ref{fvaluesremain2}.
\end{proof}


\par \medskip

\begin{claim}
\begin{align*}
 W[G_{\omega_1}] \models``y_0 & \text{ is the unique real such that }  \\& (x,y_0,m) \text{ is not coded somewhere in the $\vec{S}$-sequence.$"$ }
\end{align*}
\end{claim}
\begin{proof}[Proof of the second Claim]
We shall prove the second claim. First we show that $(x,y_0,m)$ is not coded somewhere into the $\vec{S}$-sequence. It is clear that from stage $\beta$ on, we will not code $(x,y_0,m)$ into $\vec{S}$. So the only possibility that we coded up $(x,y_0,m)$ is that there is a stage $\eta < \beta$ of our iteration $\forceP_{\omega_1}$ where we coded $(x,y_0,m)$ into $\vec{S}$. 
At stage $\eta$, we can not be in case 2, as $(x,y_0)$ and the fact that we are in case 1 at stage $\beta$, witness that we must be in case 1 at $\eta$.
So we must be in case 1, but we add a different $(x,y',m,\xi_0)$ to $I_{\eta}$, but its $<$-least name must be $<$-less than the $<$-least name for $(x,y_0,m,)$ which is a contradiction to our assumption.

In order to see that it is the unique real of the form $(x,y,m)$ which is not coded, it is sufficient to note that for every other $y \ne y_0$, $(x,y,m)$ will be coded into $\vec{S}$ by the rule (1) of our definition.

Thus Claim 2 is proved, which also finishes the proof of the Lemma under the assumptions of the first case of our Lemma.
\end{proof}


\medskip
We shall prove now that under the assumptions of the second case of our Lemma, its conclusion does hold, i.e. we need to show that if $(x,m)$ is such that for every real $y$, if $\xi < \alpha_0$ then  $(x,y,m,\xi) \notin I$, then
$W[G_{\omega_1}] \models `` \text{The $x$-section of $A_m$ is empty}".$

But under these assumptions, whenever we are at a stage $\beta$ such that there is a $y$ such that $F(\beta)=(x,y,m)$, then case 2 of the definition of $\forceP_{\omega_1}$ must apply. But for every such $y$, at stage $\beta$, we ensure with an $\alpha_0$-allowable forcing that $W[G_{\beta+1}] \models (x,y) \notin A_m$. By upwards absoluteness of $\Sigma^1_3$-formulas  we obtain  in the end \[W
[G_{\omega_1}] \models \lnot \exists y ( (x,y) \in A_m). \]
This finishes the poof of the Lemma.


\end{proof}
\begin{corollary}
In $W[G_{\omega_1}]$ the $\Pi^1_3$-uniformization property is true.
For $A_m$ an arbitrary $\Pi^1_3$-set, we get that
\begin{itemize}
\item[] $y=f(m,x)$ 
\item[]\qquad if and only if
\item[] $(x,y) \in A_m$ and $\lnot \exists r ( \forall M (M$ is countable and transitive and $ M \models \ZFP+`` \text{$\aleph_2$ exists}" $  and  $\omega_1^M=(\omega_1^L)^M $ and $ r, (x,y,m) \in M  \rightarrow M \models \varphi((x,y,m)) ).$
\end{itemize}

\end{corollary}
\begin{proof}
It suffices to note that the formula on the right and side is indeed $\Pi^1_3$. This is clear as it is of the form $\Pi^1_3 \land \lnot \Sigma^1_3$.
\end{proof}

Finally, the $\Delta^1_3$ well-order of reals follows form the fact that every pair of reals $(z_0,z_1)$ of reals got listed by $F$ and by case (3), we either coded $(z_0,z_1)$ or $(z_1,z_0)$ into $\vec{S}$. 
Thus we define in $W[G_{\omega_1}]$:
\[ z_0 < z_1 \Leftrightarrow \text{ $(z_0,z_1)$ is coded into $\vec{S}$ } \] 
the latter being a $\Sigma^1_3 (z_0,z_1)$-statement, which is what we want.

\section{Lifting the argument to $M_n$} 
We end this article with the note that our arguments presented here can be lifted to the canonical inner models with Woodin cardinals along the lines of \cite{Ho2}. Modulo the many technicalities dealt with there, one obtains 
\begin{theorem}
Let \( M_n \) be the canonical inner model with \( n \) Woodin cardinals. Then there exists a generic extension of \( M_n \) in which there is a \( \Delta^1_{n+3} \)-definable well-order of the reals, \( \mathsf{CH} \) holds, and the \( \Pi^1_{n+3} \)-uniformization property holds.
\end{theorem}





\begin{thebibliography}{12}

\bibitem{Abraham}
U. Abraham \textit{ Proper Forcing}, Handbook of Set Theory Vol.1. Springer

\bibitem{Addison}  J. Addison \textit{Some consequences of the axiom of constructibility}, Fundamenta Mathematica, vol. 46 (1959), pp. 337–357.


\bibitem{BHK} J. Baumgartner, L. Harrington and E. Kleinberg \textit{Adding a closed unbounded set.}  Journal of Symbolic Logic, 41(2), pp. 481-482, 1976.


\bibitem{David}
R. David \textit{A very absolute $\Pi^1_2$-real singleton}. Annals of Mathematical Logic 23, pp. 101-120, 1982.

\bibitem{FS}
S. D. Friedman and D. Schrittesser \textit{Projective Measure without Projective Baire}.
Memoirs of the American Mathematical Society vol. 267, 1298. 2020.

\bibitem{FH}
G. Fuchs and J. Hamkins \textit{Degrees of rigidity for Souslin Trees.}
Journal of Symbolic Logic 74(2), pp. 423-454, 2009.

\bibitem{Goldstern} M. Goldstern \textit{A Taste of Proper Forcing.} Di Prisco, Carlos Augusto (ed.) et al., Set theory: techniques and applications. Proceedings of the conferences, Curaçao, Netherlands Antilles, June 26–30, 1995 and Barcelona, Spain, June 10–14, 1996. Dordrecht: Kluwer Academic Publishers. 71-82, 1998.


    
\bibitem{Ho} S. Hoffelner \textit{$\hbox{NS}_{\omega_1}$ $\Delta_1$-definable and saturated.}  Journal of Symbolic Logic 86 (1), pp. 25 - 59, 2021.

\bibitem{Ho1} S. Hoffelner \textit{Forcing the $\Sigma^1_3$-separation property}. Journal of Mathematical Logic 22, No. 2, 2022.



\bibitem{Ho3} S. Hoffelner \textit{Forcing the $\Pi^1_3$-reduction property and a failure of $\Pi^1_3$-uniformization property},
Annals of Pure and Applied Logic,
Volume 174, Issue 8,
2023.


\bibitem{Ho2} S. Hoffelner \textit{Forcing the $\Pi^1_n$-uniformization property}, submitted. Prerprint available on  Arxiv.

\bibitem{Ho4} S. Hoffelner \textit{A Failure of $\Pi^1_{n+3}$-Reduction in the Presence of $\Sigma^1_{n+3}$-Separation}, submitted. Preprint available on Arxiv 

\bibitem{Ho5} S. Hoffelner \textit{The Global $\Sigma^1_{n+2}$-Uniformization Property and $\mathsf{BPFA}$}. Accepted at Advances in Mathematics.



\bibitem{Jech} T. Jech \textit{Set Theory. Third Millenium Edition.} Springer 2006.

\bibitem{JensenSolovay}
R. Jensen and R. Solovay \textit{Some Applications of Almost Disjoint Sets.}
Studies in Logic and the Foundations of Mathematics
Volume 59, pp. 84-104, 1970.



\bibitem{Kechris} 
A. Kechris 
\textit{Classical Descriptive Set Theory}.  
Springer 1995.




 
 \bibitem{MS} D. Martin and J. Steel \textit{A Proof of Projective Determinacy.} Journal of the American Mathematical Society (2), pp.71-125, 1989.
 
 \bibitem{Miyamoto2} T. Miyamoto {\em $\omega_1$-Suslin trees under countable support iterations.} Fundamenta Mathematicae,
vol. 143 (1993), pp. 257–261.



\bibitem{Moschovakis}
Y. Moschovakis \textit{Descriptive Set Theory.} Mathematical Surveys and Monographs 155, AMS.


\end{thebibliography}
\end{document} 


